\documentclass[11pt,a4paper]{article}

\usepackage[utf8]{inputenc}
\usepackage[T1]{fontenc}
\usepackage{amsmath,amssymb,mathtools,mathrsfs,amsthm,bm}
\usepackage{natbib}

\newtheorem{theorem}{Theorem}[section]
\newtheorem{proposition}[theorem]{Proposition}
\newtheorem{problem}[theorem]{Open problem}
\usepackage[colorlinks=true,linkcolor=blue,citecolor=blue,urlcolor=blue]{hyperref}

\newcommand{\I}{\mathscr{I}}
\newcommand{\scri}{\mathscr{I}}
\newcommand{\cT}{\mathcal{T}}
\newcommand{\cC}{\mathcal{C}}
\newcommand{\cF}{\mathcal{F}}
\newcommand{\cP}{\mathcal{P}}
\newcommand{\cE}{\mathcal{E}}
\newcommand{\cB}{\mathcal{B}}
\newcommand{\cH}{\mathcal{H}}
\newcommand{\cN}{\mathcal{N}}
\newcommand{\mathpounds}{\mathcal{L}}
\newcommand{\cR}{\mathcal{R}}
\newcommand{\cL}{\mathcal{L}}
\newcommand{\STF}{\mathrm{STF}}
\newcommand{\Div}{\operatorname{Div}}
\newcommand{\tr}{\operatorname{tr}}
\newcommand{\diff}{\mathrm{d}}

\newcommand{\statusL}{\textsf{[L]}}
\newcommand{\statusD}{\textsf{[D]}}
\newcommand{\statusF}{\textsf{[F]}}
\newcommand{\statusO}{\textsf{[O]}}

\title{Towards a Bondi--Sachs--Penrose Formalism for \(\Lambda>0\):\\
Radiation Detectors, Flux Laws, and Obstructions to News}

\author{Working manuscript}
\date{25 November 2026}

\begin{document}
\maketitle

\begin{abstract}
We assess whether a complete Bondi--Sachs--Penrose package exists for isolated systems with positive cosmological constant.  It does not: the boundary \(\scri^\pm\) is spacelike, there is no universally accepted intrinsic news tensor for generic data, no universal no-incoming-radiation condition, no general mass-positivity theorem, no BMS analogue or peeling theorem, and no universal \(\Lambda\to0\) boundary limit.  We collect the partial structures that are available: the AldS conformal--Bondi charge and flux calculus, the asymptotic super-Poynting radiation detector, the cut-wise news proposal, conditional outgoing-sector positivity, linearized de Sitter tests, exact-solution checks, and controlled flat limits.  In the no-incoming sector the rescaled supermomentum limit forces either vanishing flat Bondi news or the directional degeneracy \(B=1\), in which the incoming cut-null direction collapses.  The tensor sector of the first-jet news then has a universal \(N\)-renormalized flat limit with square determined by the outgoing rescaled supermomentum, while the full first-jet limit depends on the unresolved degeneracy rate of the incoming direction.  Several new consistency results are derived: a scaling obstruction to representing all AldS charge fluxes by one \(\xi\)-independent news square; a linear-gauge obstruction showing that a Bondi shear coefficient is fixed when \(\Lambda\neq0\); a precise harmonic characterization and round-sphere inversion of the cut-news source \(X_A\), together with a sharp \(L^2\) bound for the conditional rank-two news; finite-dimensional and representation-theoretic no-go theorems for pointwise and nonlocal linear news-square ans{\"a}tze; and a sign-flip obstruction showing that the full linearized de Sitter energy flux cannot be a nonnegative news square.  At the linearized level we derive a Hamiltonian news current whose pairings with de Sitter symmetry generators reproduce the symplectic flux, and hence the linearized AldS charge flux, with exact \(\xi\)-linearity, gauge invariance and a controlled flat limit.  We prove that the simpler volume flux law is degenerate on trace-free stress and conformal Killing symmetries, and that the radiative part of the linearized Hamiltonian news is the detector pair in the no-incoming sector.  In the strongly equipped no-incoming sector we unify the zero sets of the detector pair, radiant news, outgoing radiant supermomentum and no-radiation criterion, but we also prove a spectral norm-mismatch: the rank-two radiant news has the same zero set while being an elliptic smoothing that loses arbitrarily large high-harmonic detector weight.  Passing to the first jet of radiant news removes this mismatch: one equipment derivative recovers the tensor detector, inverse divergence recovers the vector detector, and the resulting rank-two tensor pair carries both the FAS radiant transport square and the exact detector square.  FAS VI.123 further rewrites the detector square as a local first-derivative quadratic form of the radiant news tensor itself, shifting the remaining task to coupling that news quadratic form to the AldS/symplectic flux.  The construction survives the full generalized FAS transport family: a \((\delta,\gamma)\)-rotation inversion renormalizes the first component while preserving its norm and zero set.  We further prove a charge-blindness result: nonzero linearized news can have all seven de Sitter Hamiltonian charge fluxes zero, so vanishing charge flux is strictly weaker than absence of radiation.  A conformal-weight analysis shows that the two first-jet components have weights \(\omega^{-1}\) and \(\omega^{-2}\); consequently any rank-two tensor-pair realization requires a nonlocal operator of weight \(+1\), and no natural local tensor differential operator can supply it.  We formulate a conditional symplectic dual-news representation in which the flux pairs the news with a generator-dependent dual field, clarifying why a universal news square is impossible; in the linearized boundary theory this dual field is further identified as the local adjoint of the symmetry operator on the Weyl electric field, with its radiative projection giving the canonical Riesz representative.  We formulate a retarded adjoint-factorization criterion for turning this pairing into a positive square, realize it in the ABK quadrupole sector, and normalize its radial \(\mathrm{tt}\)-projected seed so that its flat flux coincides with the standard Bondi mass-loss formula.  The remaining flat-limit coefficient \(\beta_+\) is identified exactly as the growth rate of the vector detector along the collapsing incoming cut-normal direction; bounded detector data force \(\beta_+\to0\).  In the controlled outgoing sector the FAS flat-limit relations force either vanishing flat Bondi news or the directional degeneracy \(B=1\), in which the incoming cut-null direction collapses.  We also identify the sharpest available candidate news object: in a strongly equipped, no-incoming outgoing sector the radiant news tensor is \(n^+_{AB}=2V_{AB}\), while the detector square is represented exactly by an \(\mathrm{STF}\oplus\mathrm{vector}\) Cotton pair and, on a round cut, by a minimal pair of STF tensors obtained through an \(L^2\)-isometric vector-to-tensor map.  On arbitrary \(S^2\) cuts a canonical polar decomposition extends this isometry, and a weighted conformal-Killing construction supplies the required conformally covariant weight-\(+1\) intertwiner, removing both the round-cut and covariance restrictions in the controlled sector; it is in fact the unique isometric polar factor of the local conformally covariant weighted conformal Killing operator.  Neither construction solves the unrestricted news problem.  The conditional logical core has been formalized in Lean without asserting the open physics theorem.
\end{abstract}

\tableofcontents

\section{Introduction and status labels}

The Bondi--Sachs--Penrose (BSP) formalism in asymptotically flat spacetime is a package rather than a single statement.  It contains: null infinity \(\scri^\pm\), a characteristic initial-value formulation, an intrinsic Bondi news tensor \(N_{AB}\), a no-incoming-radiation condition, Bondi mass loss and positivity, the BMS group, Nelson--Penrose peeling, and a robust scattering framework.  For \(\Lambda>0\), \(\scri^\pm\) is spacelike and this package breaks apart.  Some components survive in modified form; others have no known analogue \cite{ashtekar2017implications,poole2019bondi,poole2025charges}.

To avoid conflating different kinds of statements, we use the following labels throughout.

\begin{center}
\begin{tabular}{ll}
\statusL & rigorous statement established in the cited literature;\\[2mm]
\statusD & statement derived in this manuscript from cited equations and standard calculus;\\[2mm]
\statusF & formal or proposal-level statement, not yet a theorem of the physical theory;\\[2mm]
\statusO & open problem or target.
\end{tabular}
\end{center}

The main conclusion is:

\begin{quote}
There is currently no complete rigorous BSP package for generic isolated systems with \(\Lambda>0\).  What exists is a radiation detector, partial charge/flux laws, conditional outgoing-sector positivity, restricted monotonicity examples, and controlled flat limits for special families.
\end{quote}

\begin{table}[h]
\centering
\begin{tabular}{p{0.25\linewidth}p{0.32\linewidth}p{0.35\linewidth}}
\hline
Benchmark & \(\Lambda=0\) BSP package & \(\Lambda>0\) status\\
\hline
Infinity & null \(\scri^\pm\) & spacelike \(\scri^\pm\)\\
Gauge / well-posedness & characteristic Cauchy problem & conformal/Cauchy problem at spacelike \(\scri\); no universal Bondi gauge\\
News & intrinsic \(N_{AB}\) on 2-cuts & no universal intrinsic news; detector and conditional radiant news \(n^+_{AB}=2V_{AB}\) in strongly equipped outgoing sector\\
No incoming radiation & well-defined at \(\scri^-\) & no universal incoming/outgoing split; all radiation crosses spacelike \(\scri\)\\
Charges / mass loss & Bondi four-momentum, mass loss, positivity & AldS charges and flux laws; general positivity open\\
Symmetries & BMS & diffusion/conformal proposals; no universal BMS analogue\\
Peeling & NP peeling & no general peeling analogue\\
\(\Lambda\to0\) & built in & controlled retarded families work; universal boundary limit singular\\
\hline
\end{tabular}
\caption{Benchmark used throughout.}
\end{table}

\section{Conformal--Bondi data and AldS charge fluxes}

\subsection{Conventions}

Let \(\Lambda=3/\ell^2>0\) and set \(H=\sqrt{\Lambda/3}=1/\ell\).  We write a four-dimensional Fefferman--Graham (FG) expansion near spacelike infinity as
\begin{equation}
  g_{ab}
  =
  \frac{\ell^2}{\rho^2}
  \left[
    -\diff \rho^2
    +
    \left(
      h_{ij}+\rho^2 h^{(2)}_{ij}+\rho^3 h^{(3)}_{ij}+\cdots
    \right)\diff x^i \diff x^j
  \right].
\end{equation}
The induced metric on \(\scri\) is \(h_{ij}=g^{(0)}_{ij}\).  The holographic stress is, up to conventional powers of \(\ell\),
\begin{equation}
  \cT_{ij}
  =
  -\frac{3\ell^2}{2\kappa^2}h^{(3)}_{ij},
  \qquad \kappa^2=8\pi G,
\end{equation}
and is symmetric, trace-free, and divergence-free with respect to \(h_{ij}\) after the usual counterterm choice \cite{poole2025charges}.  Different authors absorb factors of \(\ell\) into the definition of the conformal representative; all zero sets and structural identities below are unaffected.

\subsection{Free data and conserved charges}

\statusL \textsl{Friedrich-type asymptotic data} at \(\scri\) consist of a Riemannian three-metric \(h_{ij}\) and a symmetric trace-free divergence-free tensor \(\cT_{ij}\).  These are the analogue of the characteristic free data at null infinity, but they live on a spacelike three-manifold rather than on a null hypersurface.

For a boundary vector field \(\xi^i\) and a two-surface \(S\subset\scri\) with unit normal \(n_i\) inside \(\scri\), define
\begin{equation}
  Q_\xi[S]
  =
  \int_S \cT^{ij}\xi_i n_j\,\diff S .
\end{equation}

\begin{proposition}[\statusD: conservation for conformal Killing fields]
If \(D_i\cT^{ij}=0\), \(\cT^i{}_i=0\), and \(\xi\) is a conformal Killing vector of \(h\), then \(Q_\xi[S]\) is independent of the cut \(S\).
\end{proposition}
\begin{proof}
Let \(j^i=\cT^{ij}\xi_j\).  Then
\begin{equation}
  D_i j^i
  =
  (D_i\cT^{ij})\xi_j+\cT^{ij}D_i\xi_j
  =
  \cT^{ij}D_{(i}\xi_{j)} .
\end{equation}
For a conformal Killing vector, \(D_{(i}\xi_{j)}=\frac13(D_k\xi^k)h_{ij}\), so the trace-free tensor \(\cT^{ij}\) annihilates this strain.  Hence \(D_i j^i=0\), and Stokes' theorem gives cut-independence.
\end{proof}

For a region \(B_{12}\subset\scri\) bounded by cuts \(S_1,S_2\), the flux law is
\begin{equation}\label{eq:aldS-flux}
  Q_\xi[S_2]-Q_\xi[S_1]
  =
  -\int_{B_{12}} \cT^{ij}D_i\xi_j\,\diff\mu .
\end{equation}
This is the AldS analogue of the Bondi balance law \cite{poole2025charges}.  Equivalently, since \(\cT\) is trace-free, only the trace-free part of \(D_i\xi_j\) contributes.

\subsection{Bondi-gauge obstruction to a naive news field}

\statusL In Bondi gauge with \(\Lambda\neq0\), the order-\(r^{-1}\) angular metric coefficient \(\gamma_1\) is fixed rather than free.  In the conventions of Poole--Skenderis--Taylor \cite{poole2019bondi},
\begin{equation}\label{eq:gamma1-obstruction}
  \Lambda\gamma_1
  =
  -\frac32 e^{-2\beta_0}
  \left(
    \cot\theta\,U_0
    -U_{0,\theta}
    -2U_0\gamma_{0,\theta}
    -2\gamma_{0,u}
  \right),
\end{equation}
with signs and normalizations depending on conventions.  Thus \(\gamma_1\) is algebraically determined by \((\gamma_0,\beta_0,U_0)\), unlike the asymptotically flat case where the corresponding shear coefficient is the free Bondi news candidate.

\begin{proposition}[\statusD: flat-limit scaling requirement]
For \(\gamma_1\) to remain finite as \(\Lambda\to0\) in the Bondi gauge scheme of \eqref{eq:gamma1-obstruction}, the bracket on the right-hand side must satisfy
\begin{equation}
  \cot\theta\,U_0-U_{0,\theta}-2U_0\gamma_{0,\theta}-2\gamma_{0,u}
  =
  O(\Lambda).
\end{equation}
Without this scaling the flat limit is singular at the level of the Bondi expansion.
\end{proposition}

This is a necessary condition within that gauge scheme.  It is not by itself a universal \(\Lambda\to0\) theorem.

\section{The asymptotic super-Poynting radiation detector}

\subsection{Definitions and the \(\Lambda>0\) theorem}

Let \(\Omega\) be the defining function of \(\scri\), normalized so that
\begin{equation}
  N_\alpha=\partial_\alpha\Omega,
  \qquad
  N^\alpha N_\alpha=-\frac{\Lambda}{3}.
\end{equation}
For \(\Lambda>0\), the unit timelike normal to spacelike \(\scri\) is
\begin{equation}
  n_\alpha=\sqrt{\frac{3}{\Lambda}}\,N_\alpha .
\end{equation}
Let \(d_{\alpha\beta\gamma\delta}\) denote the rescaled Weyl tensor.  Define the electric and magnetic parts relative to \(n\):
\begin{equation}
  D_{ab}=n^\mu n^\nu d_{a\mu b\nu},
  \qquad
  C_{ab}=n^\mu n^\nu\,{}^\star d_{a\mu b\nu}.
\end{equation}
Here \(D_{ab}\) is proportional to the holographic stress \(\cT_{ab}\), whereas \(C_{ab}\) is proportional to the Cotton--York tensor \(Y_{ab}\) of \((\scri,h)\); e.g. \(C_{ab}=-3Y_{ab}/\Lambda\) in a common normalization \cite{fabSen2025characterization}.  Nonzero constants are irrelevant for zero sets.

The canonical asymptotic super-Poynting vector is, up to a nonzero \(\Lambda\)-dependent factor,
\begin{equation}\label{eq:p-detector}
  p_i
  =
  2\epsilon_i{}^{jk}C_{j\ell}D^\ell{}_k
  =
  \epsilon_i{}^{jk}[C,D]_{jk},
  \qquad [C,D]=CD-DC .
\end{equation}

\begin{theorem}[\statusL: Senovilla; Fern{\'a}ndez-{\'A}lvarez--Senovilla \cite{senovilla2022radiation,fabSen2020radiation,fabSen2025characterization}]
For \(\Lambda>0\), absence of radiation on an open subset of \(\scri\) is equivalent to
\begin{equation}
  p_i=0
  \iff
  [C,D]=0 .
\end{equation}
\end{theorem}

\begin{proof}[\statusD: algebraic core]
Both \(C\) and \(D\) are symmetric, so \([C,D]\) is skew.  The map
\begin{equation}
  \mathfrak{so}(3)\to\mathbb{R}^3,
  \qquad
  \omega_{ij}\mapsto \epsilon_i{}^{jk}\omega_{jk},
\end{equation}
is the Hodge dual on two-forms and is an isomorphism.  Hence its kernel is trivial, and \(\epsilon_i{}^{jk}[C,D]_{jk}=0\) iff \([C,D]=0\).
\end{proof}

The detector is a \emph{radiation detector}, not yet news.  It tells us whether radiation is present, but it does not provide a trace-free symmetric two-tensor on cuts whose square represents the AldS charge flux.

\subsection{Why \(\Lambda<0\) is different}

For \(\Lambda<0\), \(\scri\) is Lorentzian and there is no canonical timelike normal observer.  The criterion must be imposed on the asymptotic super-Poynting vectors measured by all unit timelike observers tangent to \(\scri\).  The resulting condition is not commutation but proportionality:
\begin{equation}
  \text{no radiation}
  \iff
  C_{ab}\ \text{and}\ D_{ab}\ \text{are proportional}.
\end{equation}
Explicit photon-rocket Robinson--Trautman solutions with \(\Lambda<0\) have commuting but non-proportional Cotton--York and stress tensors, and radiate nevertheless \cite{fabSen2026rockets}.  This does not contradict the \(\Lambda>0\) theorem: there the canonical timelike normal removes the observer ambiguity.

\section{Cut-wise outgoing flux and conditional positivity}

Senovilla's cut formulation decomposes the super-Poynting vector as
\begin{equation}
  p^\alpha\;\stackrel{\scri}{=}\;\bar W\,\bar n^\alpha+\bar p^a e^\alpha{}_a,
  \qquad
  \dot{\bar W}+\bar\nabla_a\bar p^a=0 ,
\end{equation}
and introduces radiant super-momenta
\begin{equation}
  Q_\pm^\alpha
  =
  -D^\alpha{}_{\mu\nu\rho}k_\pm^\mu k_\pm^\nu k_\pm^\rho,
  \qquad
  k_\pm^\mu=\bar n^\mu\pm m^\mu ,
\end{equation}
where \(m^a\) is the unit normal to the cut inside \(\scri\) \cite{senovilla2022radiation}.  The condition \(Q_-^\mu=0\) is a proposed way to isolate an outgoing sector.

Under \(Q_-^\mu=0\), the electric Weyl data satisfy
\begin{equation}
  \cF_A=\epsilon_{AB}C^B,
  \qquad
  \hat{\cF}_{AB}=\epsilon_{AD}\hat C_B{}^D ,
\end{equation}
so the stress is fixed by the Cotton tensor except for one Coulombian scalar.  Senovilla obtains
\begin{equation}\label{eq:outgoing-pm}
  p_m=\bar p_a m^a
  =
  -2\left(
    \hat C_{AB}\hat C^{AB}+C_A C^A
  \right)
  \le 0,
  \qquad
  p_m=0\Rightarrow \bar p_a=0 .
\end{equation}
Moreover, Proposition 2 of \cite{senovilla2022radiation} gives
\begin{align}
  Q_-^\mu=Q_+^\mu=0
  &\iff Q_-^\mu=0,\ p_m=0
  \iff Q_-^\mu=0,\ \bar p_a=0\\
  &\iff
  \hat{\cF}_{AB}=\hat C_{AB}=0,
  \quad
  \cF_A=C_A=0 .
\end{align}
This is the sharpest outgoing-sector rigidity statement currently available.  It is still conditional on \(Q_-^\mu=0\) and is not by itself a construction of a news tensor.

\subsection{A detector-square news pair}

\statusD There is nevertheless a minimal news \emph{object} that reproduces the normal outgoing flux exactly.  Under \(Q_-^\mu=0\), define
\begin{equation}\label{eq:detector-news-pair}
  \cN^{\mathrm{det}}_{\mathrm{out}}
  :=
  (\hat C_{AB},\,C_A)
  \in
  \mathfrak{STF}(S)\oplus T^*S ,
  \qquad
  |\cN^{\mathrm{det}}_{\mathrm{out}}|^2
  :=
  \hat C_{AB}\hat C^{AB}+C_A C^A .
\end{equation}
Then \eqref{eq:outgoing-pm} becomes
\begin{equation}
  p_m=-2|\cN^{\mathrm{det}}_{\mathrm{out}}|^2\le0,
  \qquad
  p_m=0\iff \cN^{\mathrm{det}}_{\mathrm{out}}=0 .
\end{equation}
Combining with Proposition 2 above gives, under \(Q_-^\mu=0\),
\begin{equation}
  \cN^{\mathrm{det}}_{\mathrm{out}}=0
  \iff
  Q_-^\mu=Q_+^\mu=0 .
\end{equation}
Thus \(\cN^{\mathrm{det}}_{\mathrm{out}}\) has the required sign, positivity, and zero set for the normal outgoing flux.  Its norm has conformal weight \(\omega^{-6}\), matching the density weight of \(p_m\), so the flux identity is gauge covariant.  It evades the finite-dimensional pointwise obstruction by enlarging the target bundle from the two-dimensional space of STF tensors to \(\mathfrak{STF}(S)\oplus T^*S\).  It is not a rank-two tensor on the cut, is not the Senovilla cut news \(N_{AB}=V_{AB}+M_{AB}\), and has no proven AldS charge-flux representation.  It should therefore be viewed as a detector-square representative, not as the desired universal news tensor.

\subsection{Minimal rank-two tensor-pair representation}

\statusD The representation-theoretic obstruction says that one STF tensor has too few harmonic copies, but two STF tensors have exactly enough.  There is therefore an explicit minimal rank-two realization of the detector square on a round cut.

Fix the unit round metric and restrict to the \(l\ge2\) sector.  Let
\begin{equation}
  \mathcal K v_{AB}
  :=
  D_{(A}v_{B)}-\frac12 q_{AB}D_Cv^C
\end{equation}
be the trace-free conformal Killing operator on one-forms, and define the nonlocal operator
\begin{equation}
  \mathcal J
  :=
  \left(\frac{2}{-\Delta-2}\right)^{1/2}\mathcal K .
\end{equation}
More explicitly, write \(v_A=D_Ax+\epsilon_A{}^BD_By\).  Then
\begin{equation}
  \mathcal J v_{AB}
  =
  \sqrt{\frac{2}{-\Delta-2}}\,H_{AB}(x)
  +
  \sqrt{\frac{2}{-\Delta-2}}\,G_{AB}(y),
\end{equation}
where the square root acts on each scalar harmonic.  For \(L=l(l+1)\ge6\),
\begin{equation}
  \mathcal K(D_AY_{lm})=H_{AB}(Y_{lm}),
  \qquad
  \mathcal K(\epsilon_A{}^BD_BY_{lm})=G_{AB}(Y_{lm}),
\end{equation}
and the norm identities give
\begin{equation}
  \|\mathcal K v\|_{L^2}^2
  =
  \frac{L-2}{2}\|v\|_{L^2}^2 .
\end{equation}
Hence \(\mathcal J\) is an \(L^2\)-isometry from the \(l\ge2\) one-form sector onto the STF tensor sector, with \(\mathcal Jv=0\iff v=0\).

Define
\begin{equation}\label{eq:rank2-pair-news}
  \cN^{(2)}_{\mathrm{out}}
  :=
  (\hat C_{AB},\,\mathcal J C_A)
  \in
  \mathfrak{STF}(S)\oplus\mathfrak{STF}(S) .
\end{equation}
Then, on the radiative \(l\ge2\) sector of a round cut,
\begin{equation}
  \|\cN^{(2)}_{\mathrm{out}}\|_{L^2}^2
  =
  \int_S\big(\hat C_{AB}\hat C^{AB}+C_AC^A\big)\,dS
  =
  -\frac12\int_S p_m\,dS ,
\end{equation}
and \(\cN^{(2)}_{\mathrm{out}}=0\iff \hat C_{AB}=0,\ C_A=0\).  This saturates the harmonic multiplicity count: two STF tensors give \(4\mathcal H_l\), exactly matching the detector pair.  It is the minimal rank-two tensor realization of the outgoing detector square.

This construction has clear limitations.  It is proved on a round conformal representative; its extension to arbitrary cuts requires the conformally covariant analogue of \((-\Delta-2)^{-1/2}\), whose weight and zero-mode treatment must be fixed.  The \(l=1\) vector sector is excluded because \(-\Delta-2\) has the CKV kernel there.  Finally, \(\cN^{(2)}_{\mathrm{out}}\) is a pair of rank-two tensors, not a single rank-two news tensor, and has no proved AldS charge-flux representation.

\section{What should count as news?}

\subsection{A scaling obstruction to the oversimplified target}

The simplified target often states that a news tensor \(\mathcal N_{AB}\) should satisfy
\begin{equation}
  \Delta Q_\xi=-\int \mathcal N_{AB}\mathcal N^{AB}\,\diff\mu .
\end{equation}
This cannot hold for all boundary vector fields \(\xi\) if the right-hand side is independent of \(\xi\).

\begin{proposition}[\statusD: charge-linearity obstruction]
Assume the charge flux \(F_\xi:=\Delta Q_\xi\) is homogeneous of degree one in \(\xi\), as in \eqref{eq:aldS-flux}.  Suppose a single \(\xi\)-independent density \(\sigma\ge0\) satisfies \(F_\xi=-\int\sigma\,\diff\mu\) for every \(\xi\).  Then \(F_\xi=0\) for all \(\xi\).
\end{proposition}
\begin{proof}
For any \(a\in\mathbb{R}\), \(F_{a\xi}=aF_\xi\), whereas the proposed representation gives \(F_{a\xi}=F_\xi\).  Taking \(a=2\) yields \(2F_\xi=F_\xi\), hence \(F_\xi=0\).
\end{proof}

A viable news representation must therefore be refined in one of the following ways:
\begin{enumerate}
  \item fix a particular translation-like vector field \(\xi\);
  \item allow a \(\xi\)-dependent weight, for instance \(-\int f_\xi \mathcal N^2\,\diff\mu\) with \(f_\xi\) linear in the trace-free strain \(D_{(i}\xi_{j)}\); or
  \item introduce a news object whose bilinear pairings, not one scalar square, represent the flux functional \(\xi\mapsto\Delta Q_\xi\).
\end{enumerate}
This refinement is absent from many informal statements of the problem.

\subsection{A Hamiltonian news current and exact \(\xi\)-linearity}

\statusL The third refinement is realized exactly in the linearized de Sitter theory.  Let \(\gamma_{ab}\) denote the regular rescaled metric on \(\mathscr I^+\) and \(e_{ab}\) the rescaled electric Weyl tensor.  The symplectic potential derived from the renormalized action is \cite{hoque2025conformal}
\begin{equation}
  \Theta[\mathscr I^+]
  =
  -\frac{\ell}{16\pi G}
  \int_{\mathscr I^+}|\gamma|^{1/2}e_{ab}\,\delta\gamma^{ab}\,\diff\mu .
\end{equation}
For a de Sitter isometry whose restriction \(\xi\) to \(\mathscr I^+\) is a conformal Killing field,
\begin{equation}\label{eq:hamiltonian-flux}
  \mathcal F_\xi
  =
  \int_{\mathscr I^+}\theta(\delta\gamma,\mathcal L_\xi\delta\gamma)
  =
  \frac{\ell}{16\pi G}
  \int_{\mathscr I^+}|\gamma|^{1/2}\,\delta e^{ab}\,\mathcal L_\xi\delta\gamma_{ab}\,\diff\mu .
\end{equation}
The linearized Weyl tensor on a de Sitter background is gauge invariant, and the displayed integrand has the correct conformal weight, so \(\mathcal F_\xi\) is gauge invariant \cite{hoque2025conformal}.  This is the precise linearized form of the \(\xi\)-linear flux law: the news is not one scalar square density but a phase-space pairing between \(\delta e^{ab}\) and the symmetry generator \(\mathcal L_\xi\delta\gamma_{ab}\).

\statusD Write
\begin{equation}
  h_{ab}:=(\delta\gamma_{ab})^{\mathrm{tf}},
  \qquad
  E_{ab}:=\delta e_{ab},
\end{equation}
so that \(E^a{}_a=0\) and \(D^aE_{ab}=0\).  Since \(E^{ab}\) is trace-free, \eqref{eq:hamiltonian-flux} depends only on \(h_{ab}\).  Absorbing the constant \(\ell/(16\pi G)\) into \(\varkappa\), one obtains
\begin{equation}
  \mathcal F_\xi=\varkappa\int_{\mathscr I^+}E^{ab}\mathcal L_\xi h_{ab}\,\diff\mu .
\end{equation}
\begin{proposition}[\statusD: Hamiltonian news current]
Let \(D_{(a}\xi_{b)}=\sigma\gamma_{ab}\), with \(\sigma=\frac13D_c\xi^c\), and assume either that \(\mathscr I^+\) is compact without boundary or that the cut terms are retained.  Then
\begin{equation}\label{eq:hamiltonian-news-current}
  \mathcal F_\xi=\varkappa\int_{\mathscr I^+}\xi^c J_c[h,E]\,\diff\mu ,
\end{equation}
where
\begin{align}
  J_c[h,E]
  &:=
  -h_{ab}D_cE^{ab}
  -2D_a(E^{ab}h_{cb})
  +\frac13D_c(E^{ab}h_{ab}) \notag\\
  &=
  -\frac23 h_{ab}D_cE^{ab}
  +
  E^{ab}\Big(\frac13D_ch_{ab}-2D_ah_{cb}\Big) .
\end{align}
\end{proposition}
\begin{proof}
Use
\(\mathcal L_\xi h_{ab}=\xi^cD_ch_{ab}+h_{cb}D_a\xi^c+h_{ac}D_b\xi^c\).  The first term is integrated by parts:
\[
  \int E^{ab}\xi^cD_ch_{ab}
  =
  -\int \xi^ch_{ab}D_cE^{ab}
  -\int h_{ab}E^{ab}D_c\xi^c .
\]
Symmetry of \(E_{ab}\) gives \(\int 2E^{ab}h_{cb}D_a\xi^c=-2\int\xi^cD_a(E^{ab}h_{cb})\).  Finally \(D_c\xi^c=3\sigma\), so
\[
  -\int \sigma E^{ab}h_{ab}
  =
  \frac13\int\xi^cD_c(E^{ab}h_{ab}) .
\]
Collecting terms gives \(J_c\).  The second equality uses \(D^aE_{ab}=0\).
\end{proof}

\statusL The holographic stress tensor satisfies \(\cT_{ab}=-(\ell/8\pi G)e_{ab}\) \cite{hoque2025conformal}.  Hence \(E_{ab}=-(8\pi G/\ell)\delta\cT_{ab}\), and \eqref{eq:hamiltonian-flux} becomes
\begin{equation}\label{eq:aldS-stress-flux-linear}
  \mathcal F_\xi
  =
  -\frac12
  \int_{\mathscr I^+}
  \delta\cT^{ab}\,\delta_\xi h_{ab}\,\diff\mu ,
  \qquad
  \delta_\xi h_{ab}=\mathcal L_\xi h_{ab}+2H\beta h_{ab},
\end{equation}
with \(\beta=0\) for boundary-tangent symmetries.  This is the linearization of the candidate nonlinear AldS charge integrals, not of the simplified volume integrand in \eqref{eq:aldS-flux}.  Therefore the linearized AldS news pair may be written directly in terms of the boundary data,
\begin{equation}
  \cN^{\mathrm{AldS}}_{\mathrm{lin}}
  :=
  (h_{ab},\delta\cT^{\mathrm{tf}}_{ab}) ,
\end{equation}
rather than only in terms of \((h_{ab},E_{ab})\).

\begin{proposition}[\statusD: CKV degeneracy of the simple volume flux]\label{prop:ckv-degeneracy}
Let \(\xi\) be a conformal Killing field of \((\scri,h)\), so that
\(D_{(a}\xi_{b)}=\sigma h_{ab}\) with \(\sigma=\frac13D_c\xi^c\).  If \(\cT^{ab}\) is trace-free, then
\begin{equation}
  \cT^{ab}D_a\xi_b=0
\end{equation}
pointwise, and hence \(-\int_B\cT^{ab}D_a\xi_b\,\diff\mu=0\) for every such \(\xi\).
\end{proposition}
\begin{proof}
Decompose \(D_a\xi_b=D_{(a}\xi_{b)}+D_{[a}\xi_{b]}\).  The antisymmetric part drops out because \(\cT^{ab}\) is symmetric, while the symmetric part is pure trace and drops out because \(\cT^{ab}h_{ab}=0\).
\end{proof}

\statusL Fernández--Álvarez and Senovilla obtain the same conclusion for their first-class charges \cite{fabSen2022asymptotic}.  For a first-class current \(y^a=t^{ab}\xi_b\) with \(t^{ab}\) built from \(D_{ab}\), \(C_{ab}\), or their outgoing/incoming parts, their equations (VII.98)--(VII.107) give a balance law sourced only by matter terms and Coulomb contributions; gravitational radiation does not enter.  They explicitly note that modifications of these currents are needed.  Thus any news candidate aiming to represent AldS radiation flux cannot be obtained from first-class stress/Cotton moments alone; it must involve symplectic/Hamiltonian data, second-class Bel--Robinson currents, or radiant news and its time derivative.

\statusL This degeneracy is consistent with Ashtekar--Bonga--Kesavan's observation that imposing conformal flatness of the intrinsic geometry of \(\mathscr I^+\) from the outset trivializes the flux problem \cite{ashtekar2015linear}.  The nonzero Hamiltonian flux \eqref{eq:hamiltonian-flux} survives because the candidate nonlinear charges depend on the full boundary conformal geometry; their linearization contains \(\mathcal L_\xi h_{ab}\), not merely \(\cT^{ab}D_a\xi_b\) evaluated on the background.  Thus \eqref{eq:aldS-stress-flux-linear} is the correct linearized charge-flux relation, while \eqref{eq:aldS-flux} is a special balance identity that vanishes identically on trace-free stress and conformal symmetries.

\statusL Senovilla's TT tensor \(\mathcal F_{ab}\) is the electric part of the rescaled Weyl tensor \cite{senovilla2022radiation}, so at linear order \(\delta\cT_{ab}\) is proportional to \(\mathcal F_{ab}\).  \statusD Under the no-incoming condition \(Q_-^\mu=0\),
\begin{equation}
  \mathcal F_A=\epsilon_{AB}C^B,
  \qquad
  \hat{\mathcal F}_{AB}=\epsilon_{AD}\hat C_B{}^D ,
\end{equation}
as recalled in \eqref{eq:outgoing-pm}.  Consequently the radiative cut projection of \(\cN^{\mathrm{AldS}}_{\mathrm{lin}}\) is, up to the fixed \(\epsilon\)-rotation and convention-dependent normalization, exactly the detector pair \((\hat C_{AB},C_A)\):
\begin{equation}\label{eq:linear-news-detector-bridge}
  \Pi_{\mathrm{rad}}\big(\cN^{\mathrm{AldS}}_{\mathrm{lin}}\big)
  \longleftrightarrow
  (\hat C_{AB},C_A),
  \qquad
  \big|\Pi_{\mathrm{rad}}\delta\cT\big|^2
  \propto
  \hat C_{AB}\hat C^{AB}+C_AC^A
  =
  -\tfrac12 p_m .
\end{equation}
Thus, in the linearized no-incoming sector, the object that pairs with \(\xi\) to give the AldS charge flux and the object that detects outgoing radiation have the same radiative components and the same zero set.  The remaining nonradiative component \(\mathcal F_{ab}m^am^b\) affects tangential super-Poynting components but not \(p_m\).  This is a compatibility theorem, not a nonlinear news construction: it assumes the \(Q_-^\mu=0\) sector, fixes a cut projection, and does not by itself prove positivity on the full phase space.

This current supplies the missing linearized realization of the news-as-pairing requirement.  The object
\begin{equation}\label{eq:hamiltonian-news-pair}
  \cN^{\mathrm H}:=(h_{ab},E_{ab})
\end{equation}
is a phase-space pair, and the flux is \(\mathcal F_\xi=\varkappa\langle \cN^{\mathrm H},\delta_\xi h\rangle\) with \(\delta_\xi h=\mathcal L_\xi h\) for boundary-tangent symmetries.  If the spacetime Killing field extends off \(\mathscr I^+\) with normal component \(\beta\), the boundary generator acquires the extra conformal term \(2H\beta\), giving the Ashtekar--Bonga--Kesavan flux \cite{ashtekar2015linear}
\begin{equation}\label{eq:abk-hamiltonian-flux}
  \mathcal H_T
  =
  \frac{1}{2H\kappa}
  \int_{\mathscr I^+}
  E^{ab}\big(\mathcal L_Th_{ab}+2Hh_{ab}\big)\,\diff\mu .
\end{equation}
For smooth retarded families this reduces to the positive Minkowski flux
\begin{equation}
  \mathcal H_T\longrightarrow
  \frac1{8\kappa}
  \int_{\mathscr I^+}
  \Big[
  (\partial_t\mathring h_{ab})(\partial_t\mathring h^{ab})
  +
  (D_m\mathring h_{ab})(D^m\mathring h^{ab})
  \Big]\,\diff\mu ,
\end{equation}
so the linearized Hamiltonian news has the correct controlled flat limit \cite{ashtekar2015linear}.

\statusL Ashtekar--Bonga--Kesavan also emphasize that, for their prescription of asymptotic symmetries, the difference of the candidate nonlinear charge integrals on \(\mathscr I^+\) linearizes precisely to the symplectic flux formulas \eqref{eq:abk-hamiltonian-flux} and its momentum counterparts \cite{ashtekar2015linear}.  Hence \(\cN^{\mathrm H}\) is not merely an abstract symplectic current: it represents the \emph{linearization of the AldS charge-flux law} itself.  This gives a rigorous linearized answer to the \(\xi\)-linearity part of the central problem, while leaving the nonlinear construction open.

This is significant but not a solution of the nonlinear central problem.  It proves exact \(\xi\)-linearity, gauge invariance and the controlled flat limit in the linearized sector.  It does not make \(\mathcal F_\xi\) a nonnegative square on the full phase space: the sign-flip obstruction already shows that \(E\mapsto-E\), \(B\mapsto B\) reverses the de Sitter energy while preserving the detector square.  Positivity still requires extra incoming/outgoing constraints.  It is also a flux pairing with a phase-space derivative of \(h_{ab}\), not a rank-two tensor on a cut and not a proven representation of the nonlinear AldS charge \(\cT^{ij}\).  The finite algebra is formalized by \texttt{HamiltonianNewsFlux.Model}, \texttt{flux\_smul}, and \texttt{flux\_add}.

\subsection{The cut-news proposal}

\statusF Senovilla's 2022 program proposes a cut-wise news tensor
\begin{equation}\label{eq:N-decomposition}
  N_{AB}=V_{AB}+M_{AB}
\end{equation}
on a compact spherical cut \(S\subset\scri^+\).  Let \(q_{AB}\) be the induced metric, \(m^a\) the unit normal in \(\scri\), \(\varkappa_{AB}\) the second fundamental form, and write
\begin{equation}
  \varkappa_{AB}=\Sigma_{AB}+\frac12\varkappa q_{AB},
  \qquad \Sigma^A{}_A=0 .
\end{equation}
Let \(S_{AB}\) be the Schouten tensor of the cut and let \(\rho_{AB}\) be the unique Codazzi tensor of trace equal to the Gaussian curvature \(K\).  Then
\begin{equation}
  U_{AB}
  =
  S_{AB}
  -\frac12\Sigma_A{}^D\Sigma_{BD}
  +\frac{\varkappa}{2}\Sigma_{AB}
  +\frac{\varkappa^2}{8}q_{AB},
  \qquad
  V_{AB}=U_{AB}-\rho_{AB}.
\end{equation}
The tensor \(V_{AB}\) is symmetric, trace-free, gauge invariant, and intrinsic to \((\scri,h)\) and the cut.  The unknown part \(M_{AB}\) is symmetric, trace-free and gauge invariant, and satisfies
\begin{equation}
  D_{[C}M_{A]B}
  =
  X_{CAB}
  -\frac12\sqrt{\frac{\Lambda}{3}}\epsilon_{CA}C_B
  +D_{[C}\!\left(\Sigma_{A]E}\Sigma_B{}^E\right)
  -\frac12 D_B\Sigma_{[C}{}^E\Sigma_{A]E},
\end{equation}
with
\begin{equation}
  X_{CAB}=X_{[CA]B}=2q_{B[C}X_{A]} .
\end{equation}
Equivalently,
\begin{equation}\label{eq:X-source-div}
  \frac12 D_C M^C{}_A
  =
  X_A
  +\frac12\sqrt{\frac{\Lambda}{3}}\epsilon_{AB}C^B
  -\frac38 D_A(\Sigma_{DE}\Sigma^{DE})
  +\frac12\Sigma^{CE}D_C\Sigma_{EA}.
\end{equation}
Finally,
\begin{equation}
  D_{[C}N_{A]B}=X_{CAB}.
\end{equation}
The essential open input is the choice of the one-form \(X_A\).

\subsection{Constraints on \(X_A\)}

\statusD \textsl{Gauge weight.}  Under a conformal rescaling of the cut metric, \(X_A\) must transform as
\begin{equation}
  \tilde X_A=\omega^{-2}X_A .
\end{equation}
The canonical super-Poynting vector instead has weight \(\omega^{-5}\), so a scalar built directly from \([C,\cF]\) does not automatically have the correct weight to serve as the dimensionless gauge-invariant input of the ansatz below.

\statusL \textsl{CKV integrability.}  Solvability of the Codazzi equation requires six orthogonality conditions on \(S^2\):
\begin{equation}\label{eq:CKV-integrability}
  \int_S \xi^A X_A\,\diff S=0
  \qquad
  \text{for every conformal Killing vector }\xi^A .
\end{equation}
This is eq. (77) of \cite{senovilla2022radiation}.

\statusD \textsl{Hodge content.}  On \(S^2\) there are no harmonic one-forms, so
\begin{equation}
  X_A=\epsilon_A{}^B D_B x+D_A X ,
\end{equation}
with potentials defined only up to additive constants.  Condition \eqref{eq:CKV-integrability} removes the \(s=1\) spherical-harmonic parts of \(x\) and \(X\); the \(s=0\) parts are gauge.  Thus admissible sources have only \(s\ge2\) harmonic content.

\statusL \textsl{Suggested ansatz.}  Senovilla suggests
\begin{equation}\label{eq:X-ansatz}
  X_A=\Delta f\,D_A f
\end{equation}
with gauge-invariant dimensionless \(f\).  Lemma 4 of \cite{senovilla2022radiation} gives \(\int_S \Delta f\,\mathpounds_\xi f=0\) for every CKV \(\xi\), so the ansatz automatically satisfies integrability.  On a connected \(S^2\), \(X_A=0\) iff \(f\) is constant.

No canonical \(f\) is known.  In the outgoing sector one might try to build \(f\) from Cotton and \(\cF\) potentials, but the potentials in the decomposition of STF tensors on \(S^2\) are ambiguous by \(s=0,1\) pieces, and the paper supplies no formula linking \(X_A\) to \(\bar p_a\) or \(Q_-^\mu\).  The bridge is only conceptual: in the unrestricted problem \(V_{AB}\) is intrinsic geometry, so \(M_{AB}\), hence \(X_A\), must carry the \(\cF_{ab}\)-dependence absent from \(V_{AB}\).  The strongly equipped outgoing sector below is an important exception because the no-incoming condition itself determines the radiative electric data from intrinsic Cotton data.

\subsection{A conditional rank-two news object: strongly equipped radiant news}

\statusL Fern{\'a}ndez-{\'A}lvarez and Senovilla define an \emph{equipped} \(\scri^+\) by a chosen unit tangent field \(m^a\); it is \emph{strictly equipped} when \(m^a\) is surface-orthogonal and defines a foliation by cuts, and \emph{strongly equipped} when the cuts are additionally umbilical \cite{fabSen2022asymptotic}.  On each cut they introduce radiant news tensors
\begin{equation}
  n^\pm_{AB}=V_{AB}+X^\pm_{AB},
\end{equation}
where \(X^\pm_{AB}\) are second components constrained by the incoming/outgoing radiant-super-momentum equations.  For strictly equipped \(\scri^+\) with \(S^2\)-topological leaves,
\begin{equation}
  n^+_{AB}=0\iff Z^+=0,
  \qquad
  n^-_{AB}=0\iff Z^-=0 ,
\end{equation}
where \(Z^\pm\) are the transverse radiant-super-momentum components.  If \(\scri^+\) is strongly equipped, has \(S^2\) leaves, and satisfies the no-incoming condition \(Q_-^\mu=0\), then the outgoing second component can be chosen as \(X^+_{AB}=V_{AB}\), so that
\begin{equation}\label{eq:strong-news}
  n^+_{AB}=2V_{AB}.
\end{equation}
This is the sharpest rank-two news candidate obtained from the available literature.

The zero-set statement is stronger than the outgoing radiant supermomentum criterion.  Theorem 2 of \cite{fabSen2022asymptotic} states that, under criterion 2 (no incoming radiation), strict equipment with \(S^2\) leaves, and the two radiant-news solvability equations,
\begin{equation}
  n^+_{AB}=0
  \iff
  P^a=0
  \iff
  \text{no radiation at }\scri^+ ,
\end{equation}
where \(P^a\) is the canonical asymptotic super-Poynting vector.  Corollary VI.8 specializes this to strong equipment and gives
\begin{equation}\label{eq:V-zero-P-zero}
  V_{AB}=0
  \iff
  P^a=0
  \iff
  \text{no radiation at }\scri^+ .
\end{equation}
Thus \(V_{AB}\), and equivalently \(n^+_{AB}=2V_{AB}\), is a genuine rank-two news tensor inside this controlled sector: it is symmetric, trace-free, gauge invariant, intrinsic to \((\scri,h)\), and has exactly the radiation-detector zero set.

\statusF If Senovilla's cut-news tensor is identified with the Fernández--Álvarez--Senovilla outgoing radiant news in this sector, then \(X^+_{AB}=V_{AB}\) corresponds to \(M_{AB}=V_{AB}\).  Under this natural but unproved bridge, umbilicity gives \(\Sigma_{AB}=0\), so the Codazzi-source equation \eqref{eq:X-source-div} determines the missing one-form explicitly:
\begin{equation}\label{eq:X-strong-out}
  X_A
  =
  \frac12 D_C V^C{}_A
  -
  \frac12\sqrt{\frac{\Lambda}{3}}\,\epsilon_{AB}C^B .
\end{equation}
Under criterion 2 the Cotton vector \(C^B\) is intrinsic data, so \(X_A\) is intrinsic in this sector.  This supplies the previously missing Codazzi source for the cut-news ansatz, but only in the strongly equipped outgoing sector and only modulo the cross-formalism identification.  The reverse implication \(X_A=0\Rightarrow V_{AB}=0\) is not automatic: \eqref{eq:X-strong-out} is a first-order relation between \(V\) and \(C\), not by itself a rigidity statement.

\statusD A more direct nonlocal formula is available entirely within the radiant-news formalism.  Lemma V.3, equations (V.70) and (V.71), and the proof of (V.88) in \cite{fabSen2022asymptotic} give
\begin{equation}\label{eq:div-V-strong-out}
  D^B V_{AB}
  =
  -N\,\epsilon_{AB}C^{+B} .
\end{equation}
Thus \(V_{AB}\) is the inverse-divergence of the outgoing Cotton one-form.  On a unit round cut, write the source \(S_A:=-N\epsilon_{AB}C^{+B}\) in harmonics as
\begin{equation}
  S_A
  =
  \sum_{l\ge2,m}
  \Big(
  \sigma_{lm}D_AY_{lm}
  +
  \tau_{lm}\epsilon_A{}^BD_BY_{lm}
  \Big),
\end{equation}
with \(l\le1\) removed by the CKV solvability conditions and gauge freedom.  Using the divergences of the STF harmonics given above,
\begin{equation}
  D^A H_{AB}(Y_{lm})
  =
  \left(1-\frac{l(l+1)}2\right)D_BY_{lm},
  \qquad
  D^A G_{AB}(Y_{lm})
  =
  \frac{1-l(l+1)}2\,\epsilon_B{}^CD_CY_{lm},
\end{equation}
one obtains
\begin{equation}\label{eq:V-round-inversion}
  V_{AB}
  =
  \sum_{l\ge2,m}
  \left[
  \frac{\sigma_{lm}}{1-l(l+1)/2}\,H_{AB}(Y_{lm})
  +
  \frac{2\tau_{lm}}{1-l(l+1)}\,G_{AB}(Y_{lm})
  \right],
  \qquad
  n^+_{AB}=2V_{AB}.
\end{equation}
This is an explicit nonlocal rank-two news object in the strongly equipped outgoing sector.  It shows that even the genuine \(n^+=2V\) news is inverse-elliptic on each cut rather than an algebraic shear derivative.

\statusD The inversion also gives a sharp cut-wise elliptic estimate.  Let \(L=l(l+1)\).  Bochner's identity on the unit round sphere, the relation \(G=\frac12 JH\) for the Hodge dual \(J\) of an STF two-tensor, and orthogonality of gradient and curl vector harmonics give
\begin{equation}
  \int_S |H_{AB}(Y_{lm})|^2\,dS
  =
  \int_S |G_{AB}(Y_{lm})|^2\,dS
  =
  \frac{L(L-2)}2\int_S |Y_{lm}|^2\,dS,
  \qquad
  \int_S |D_AY_{lm}|^2\,dS
  =
  \int_S |\epsilon_A{}^BD_BY_{lm}|^2\,dS
  =
  L\int_S |Y_{lm}|^2\,dS .
\end{equation}
Substituting the multipliers in \eqref{eq:V-round-inversion} therefore yields
\begin{equation}\label{eq:V-round-L2}
  \int_S |V_{AB}|^2\,dS
  =
  \sum_{l\ge2,m}
  \left[
  \frac{2}{L-2}\,|S^{\mathrm{grad}}_{lm}|^2
  +
  \frac{2(L-2)}{(L-1)^2}\,|S^{\mathrm{curl}}_{lm}|^2
  \right],
\end{equation}
where \(S^{\mathrm{grad}}_{lm}=\sigma_{lm}D_AY_{lm}\) and \(S^{\mathrm{curl}}_{lm}=\tau_{lm}\epsilon_A{}^BD_BY_{lm}\).  Since \(L\ge6\),
\begin{equation}\label{eq:V-round-bound}
  \int_S |V_{AB}|^2\,dS
  \;\le\;
  \frac12\int_S |S_A|^2\,dS,
  \qquad
  \int_S |n^+_{AB}|^2\,dS
  \;\le\;
  2\int_S |S_A|^2\,dS ,
\end{equation}
with equality in the first inequality only for a pure \(l=2\) gradient source.  This is a quantitative inverse-elliptic smoothing statement.  It does not convert the Cotton detector into an AldS flux news.

\statusD The same computation gives a sharper magnitude obstruction.  For a pure gradient or curl harmonic source,
\begin{equation}
  \frac{\int_S |V_{AB}|^2\,dS}{\int_S |S_A|^2\,dS}
  =
  \begin{cases}
    \dfrac{2}{L-2}, & S_A \text{ pure gradient},\\[1ex]
    \dfrac{2(L-2)}{(L-1)^2}, & S_A \text{ pure curl}.
  \end{cases}
\end{equation}
Both ratios tend to zero as \(L=l(l+1)\to\infty\).  Hence there is no constant \(c>0\) such that
\begin{equation}
  c\int_S |S_A|^2\,dS
  \;\le\;
  \int_S |V_{AB}|^2\,dS
\end{equation}
for all sources.  Since \(S_A=-N\epsilon_{AB}C^{+B}\) is the outgoing Cotton source of the radiant supermomentum, the rank-two news \(n^+=2V\) has the same zero set as the Cotton detector but a strictly weaker high-harmonic norm: it is an order-\((-1)\) elliptic smoothing of the detector, not a norm-equivalent representative of its square.  This explains precisely why the unified zero set \eqref{eq:outgoing-zero-unification} does not by itself give a central news-square theorem.  A nonlinear object reproducing the detector magnitude must retain the missing high-harmonic spectral weight, for example through a derivative, a two-tensor pair, or a bilinear flux pairing rather than through \(\int |n^+|^2\) alone.

\begin{proposition}[\statusF: spectral norm mismatch]
Let \(D[S]\) be the Cotton-detector quadratic norm and \(N[S]:=\int_S|V|^2\) the rank-two news norm.  The pair \((D,N)\) has spectral mismatch if \(N\le C D\) for some \(C>0\) but no \(c>0\) satisfies \(cD\le N\) for all sources.  The round-cut identities above show that \((D,N)\) has spectral mismatch on the strongly equipped outgoing sector.
\end{proposition}

\begin{proposition}[\statusD: conditional radiant-news object]
Under the hypotheses of Theorem 2 and Corollary VI.8 of \cite{fabSen2022asymptotic}, \(n^+_{AB}=2V_{AB}\) is a symmetric trace-free gauge-invariant rank-two tensor on each cut and
\begin{equation}
  n^+_{AB}=0
  \iff
  V_{AB}=0
  \iff
  P^a=0 .
\end{equation}
If in addition \(p_m=0\), then Proposition 2 of \cite{senovilla2022radiation} gives \(Q_+^\mu=Q_-^\mu=0\), so the outgoing and incoming radiant sectors are both absent.
\end{proposition}
\begin{proof}
Symmetry, trace-freeness and gauge invariance are the content of Propositions V.1 and VI.2 of \cite{fabSen2022asymptotic}; the identity \(n^+=2V\) is Lemma VI.10 there.  The zero set follows from Theorem 2 and Corollary VI.8 of the same paper.  The last assertion follows from the equivalence \(Q_-^\mu=0,\ p_m=0\iff Q_-^\mu=Q_+^\mu=0\) in Proposition 2 of \cite{senovilla2022radiation}.
\end{proof}

\subsection{Unification of the outgoing zero sets}

\begin{theorem}[\statusD: common zero set of detector, radiant news and outgoing supermomentum]
Assume \(Q_-^\mu=0\) and the hypotheses of Proposition VI.4, Theorem 2 and Corollary VI.8 of \cite{fabSen2022asymptotic}: strict equipment, \(S^2\) leaves, criterion 2, and the radiant-news solvability equations; for the \(V_{AB}\) statement assume strong equipment.  Then
\begin{equation}\label{eq:outgoing-zero-unification}
  \cN^{\mathrm{det}}_{\mathrm{out}}=0
  \iff
  Q_+^\mu=0
  \iff
  n^+_{AB}=0
  \iff
  V_{AB}=0
  \iff
  P^a=0
  \iff
  \text{no radiation at }\scri^+ .
\end{equation}
\end{theorem}
\begin{proof}
Under \(Q_-^\mu=0\), Senovilla's relations \(\hat\cF_{AB}=\epsilon_{AD}\hat C_B{}^D\) and \(\cF_A=\epsilon_{AB}C^B\) show that \(\cN^{\mathrm{det}}_{\mathrm{out}}=0\) is equivalent to \(\hat\cF_{AB}=0=\cF_A\).  Proposition 2 of \cite{senovilla2022radiation} then gives \(Q_+^\mu=0\).  Proposition VI.4 of \cite{fabSen2022asymptotic} gives \(Q_+^\mu=0\iff n^+_{AB}=0\).  Lemma VI.10 gives \(n^+_{AB}=2V_{AB}\), Proposition VI.4 and Theorem 2 give the equivalence with \(P^a=0\), and Corollary VI.8 gives the strong-equipment statement for \(V_{AB}\).
\end{proof}

This is the sharpest conditional unification currently available.  It says that the detector pair, the rank-two radiant news, and the outgoing radiant supermomentum all vanish on exactly the same physical sector.  It does not yet identify their \emph{magnitudes} with one AldS charge-flux functional: \(\cN^{\mathrm{det}}_{\mathrm{out}}\) has the exact detector square, \(n^+_{AB}=2V_{AB}\) has the FAS transport square through \(\dot n^+\), and the linearized Hamiltonian pair has the symplectic/AldS flux pairing.  The central theorem requires a single object retaining all three structures.

\statusL The same reference gives a transport law that clarifies what the radiant news measures.  Under strict equipment, \(S^2\) leaves, and equations (VI.86), (VI.87), (VI.111), (VI.112) of \cite{fabSen2022asymptotic}, equations (VI.106)--(VI.109) give
\begin{equation}\label{eq:radiant-transport}
  N\,\epsilon_B{}^E C^{(\pm)}_{AE}
  =
  \dot n^{(\pm)}_{AB}
  -
  \Sigma^{(\pm)}_{(AC}n^{(\pm)}_{B)}{}^C ,
  \qquad
  N^2 W^{(\pm)}
  =
  \Big(\dot n^{(\pm)}_{AB}
  -
  \Sigma^{(\pm)}_{(AC}n^{(\pm)}_{B)}{}^C\Big)
  \Big(\dot n^{(\pm)AB}
  -
  \Sigma^{(\pm)A}{}_{D}n^{(\pm)BD}\Big) .
\end{equation}
For umbilical leaves, \(\Sigma^{(\pm)}_{AB}=0\), so the superenergy density is exactly the square of the evolution derivative of the radiant news:
\begin{equation}\label{eq:radiant-flux-square}
  N^2 W^{(\pm)}=\dot n^{(\pm)}_{AB}\dot n^{(\pm)AB}.
\end{equation}
In the strongly equipped outgoing sector this becomes \(N^2W^+=4\dot V_{AB}\dot V^{AB}\).  This is a genuine flux-square representation at the level of the FAS radiant superenergy.  It is not an AldS charge-flux law: \(W^+\) is a super-Poynting density, the derivative is taken along the chosen equipment, and the existing zero-set theorem concerns \(n^+\) itself rather than \(\dot n^+\).

This is not a solution of the unrestricted news problem.  It requires a chosen equipment, umbilical cuts, \(S^2\) topology, and the no-incoming condition.  It has not been shown that \(n^+_{AB}\) reproduces all AldS charge fluxes with the correct linearity in \(\xi\), nor that it has a universal \(\Lambda\to0\) limit.  What it does provide is a genuine conditional rank-two tensor whose zero set matches the outgoing radiation detector in the sharpest controlled sector currently available.

\subsection{First-jet radiant news: restoring the detector norm}

\statusD The spectral mismatch of \(n^+=2V\) is removed by passing from the news tensor to its first jet along the equipment.  Keep strict equipment, \(S^2\) leaves, criterion 2 \(Q_-^\mu=0\), and the radiant solvability equations; for the divergence step assume strong equipment.  Let \(J\) be the Hodge operator on STF two-tensors and \(j\) the Hodge operator on one-forms, with \(J^2=j^2=-1\).  Define the flux part of the radiant jet by
\begin{equation}\label{eq:first-jet-flux}
  \Phi^+_{AB}
  \;:=\;
  \frac{1}{2N}
  \Big(
  \dot n^+_{AB}
  -
  \Sigma^+_{(AC}n^+_{B)}{}^C
  \Big) .
\end{equation}
The FAS transport law \eqref{eq:radiant-transport} gives
\(\Phi^+_{AB}=\tfrac12\epsilon_B{}^E\hat C^{(+)}_{AE}\), where \(\hat C^{(+)}_{AB}\) is the outgoing lightlike projection of the Cotton tensor.  Senovilla's no-incoming relations \(\hat{\cF}=J\hat C\), \(\cF_A=jC_A\), combined with properties xv) and xxi) of the FAS lightlike projections, give
\begin{equation}\label{eq:outgoing-cotton-projections}
  \hat C^{(+)}_{AB}=2\hat C_{AB},
  \qquad
  \hat C^{(+)}_A=C_A ,
\end{equation}
up to the fixed Hodge/sign convention.  Therefore
\begin{equation}\label{eq:first-jet-flux-detector}
  \Phi^+_{AB}
  =
  \epsilon_B{}^E\hat C_{AE},
  \qquad
  \Phi^+_{AB}\Phi^{+AB}
  =
  \hat C_{AB}\hat C^{AB} .
\end{equation}
Thus the first jet component is the normalized FAS flux density: its square is \(W^+/4\) and equals the tensor part of the detector square.  Explicitly, FAS VI.108 and the definition \eqref{eq:first-jet-flux} give
\begin{equation}\label{eq:W-phi-normalization}
  N^2W^+
  =
  \big(\dot n^+_{AB}-\Sigma^+_{(AC}n^+_{B)}{}^C\big)
  \big(\dot n^{+AB}-\Sigma^{+A}{}_{D}n^{+BD}\big)
  =
  4N^2\Phi_{AB}\Phi^{AB},
\end{equation}
so \(W^+=4\Phi^2=4\hat C^2\) in the outgoing no-incoming sector.  This is the normalization used throughout; the larger coefficient \(8\) that appeared in an earlier draft of the flat-scaling subsection was a double counting of the Senovilla and FAS radiant-superenergy conventions.

\statusD For the vector part, \(n^+=2V\) and the inverse-divergence law \eqref{eq:div-V-strong-out} give
\begin{equation}\label{eq:vector-from-news}
  D^B n^+_{AB}
  =
  -2N\,\epsilon_{AB}C^B ,
\end{equation}
where under \(Q_-^\mu=0\) the \(C^B\) on the right is the detector Cotton vector.  Let \(C_A[n^+]\) denote the unique \(l\ge2\) one-form solving \eqref{eq:vector-from-news}; its round-cut inversion is the one already used in \eqref{eq:V-round-inversion}.  On a round cut, apply the isometric operator \(\mathcal J\) of \eqref{eq:rank2-pair-news} and define the first-jet radiant news
\begin{equation}\label{eq:first-jet-news}
  \cN^{\mathrm{jet}}_+
  \;:=\;
  \big(\Phi^+_{AB},\,\mathcal J C_A[n^+]\big)
  \in
  \mathfrak{STF}(S)\oplus\mathfrak{STF}(S) .
\end{equation}
Then
\begin{equation}\label{eq:first-jet-square}
  \|\cN^{\mathrm{jet}}_+\|_{L^2}^2
  =
  \int_S\big(
  \hat C_{AB}\hat C^{AB}+C_AC^A
  \big)\,dS
  =
  -\frac12\int_S p_m\,dS ,
\end{equation}
and
\begin{equation}\label{eq:first-jet-zero}
  \cN^{\mathrm{jet}}_+=0
  \iff
  \hat C_{AB}=0,\ C_A=0
  \iff
  Q_-^\mu=Q_+^\mu=0
  \iff
  n^+_{AB}=0
  \iff
  \text{no radiation at }\scri^+ .
\end{equation}
The first equivalence uses \eqref{eq:first-jet-flux-detector} and the fact that \(\mathcal J\) is an isometry on the radiative \(l\ge2\) sector; the remaining equivalences are the outgoing zero-set theorem \eqref{eq:outgoing-zero-unification}.

\statusL FAS VI.123 supplies a second, news-only expression for the vector detector square.  With gauge scalars \(\beta_{\mathrm F},\lambda_{\mathrm F}\) and \(Z^+\) the FAS vector superenergy,
\begin{equation}\label{eq:FAS-divergence-superenergy}
  N^2\big(\beta_{\mathrm F}^2+\lambda_{\mathrm F}^2\big)Z^+
  =
  D_Cn^{+B C}D_Dn^{+BD} .
\end{equation}
In the controlled outgoing sector \(\beta_{\mathrm F}^2+\lambda_{\mathrm F}^2=1\), and the inverse-divergence law \eqref{eq:vector-from-news} gives \(D^Bn^+_{AB}=-2N\epsilon_{AB}C^B\), so \eqref{eq:FAS-divergence-superenergy} reduces to \(Z^+=4C^2\), in agreement with \eqref{eq:alpha-beta-W-Z}.  Combining it with \eqref{eq:W-phi-normalization} yields the manifestly news-only quadratic form
\begin{equation}\label{eq:jet-news-only-square}
  N^2\,\|\cN^{\mathrm{jet}}_+\|_{L^2}^2
  =
  \frac14\int_S
  \Big[
    \big(\dot n^+_{AB}-\Sigma^+_{(AC}n^+_{B)}{}^C\big)
    \big(\dot n^{+AB}-\Sigma^{+A}{}_{D}n^{+BD}\big)
    +
    D_Cn^{+B C}D_Dn^{+BD}
  \Big]\,dS .
\end{equation}
Thus the detector square is not merely reconstructed from the detector: in the controlled sector it is exactly a first-derivative quadratic form of the radiant news tensor.  This reframes the remaining coupling problem: one must pair this news quadratic form with the AldS/symplectic flux, not search for a different detector square.

This is the first candidate that simultaneously carries the unified outgoing zero set, the FAS radiant transport square, the exact detector square, and the minimal rank-two tensor-pair structure.  It also explains the earlier spectral mismatch: \(n^+=2V\) loses high-harmonic detector weight because the vector detector is recovered by inverse divergence, while the tensor detector is recovered by one equipment derivative; the first jet retains both operations.  The construction remains conditional and incomplete.  It uses the first jet rather than data on a single cut, it is tied to strict/strong equipment and an \(S^2\) foliation, its conformal covariance on arbitrary cuts requires the weighted intertwiner developed below and exclusion of conformal-Killing components, and it has not been paired with a general AldS/symplectic charge flux or with a universal \(\Lambda\to0\) limit.  The conditional finite algebra is formalized by \texttt{RadiantJetNews.Model}, \texttt{jet\_sq\_eq\_neg\_half\_normalFlux}, and \texttt{jet\_zero\_iff\_detector\_zero}.

\subsection{Generalized \((\delta,\gamma)\) first-jet inversion}

\statusD The umbilical normalization is not essential.  FAS allow generalized transport equations with gauge-invariant dimensionless scalars \(\delta,\gamma\) satisfying \(\delta^2+\gamma^2>0\) \cite{fabSen2022asymptotic}:
\begin{equation}\label{eq:generalized-transport}
  N\big(
    \delta\,\cC^{(+)}_{AB}
    +\gamma\,\epsilon_{(A}{}^{C}\cC^{(+)}_{B)C}
  \big)
  =
  \dot n^+_{AB}-\Sigma^+_{(AC}n^+_{B)}{}^C ,
\end{equation}
where \(\cC^{(+)}_{AB}\) is the outgoing tensor detector projection.  Let \(R\) denote the rotation \((RT)_{AB}:=\epsilon_{(A}{}^{C}T_{B)C}\) on STF tensors.  On any two-dimensional Riemannian cut, \(R^2=-1\) on the radiative STF sector and
\begin{equation}
  \big\|(\delta I-\gamma R)T\big\|^2
  =
  (\delta^2+\gamma^2)\|T\|^2 .
\end{equation}
Define the unnormalized jet again by
\begin{equation}
  \Phi_{AB}
  :=
  \frac{1}{2N}
  \big(
    \dot n^+_{AB}-\Sigma^+_{(AC}n^+_{B)}{}^C
  \big).
\end{equation}
Then \eqref{eq:generalized-transport} gives
\begin{equation}
  \Phi_{AB}
  =
  \frac12\big(\delta\,\cC^{(+)}_{AB}+\gamma(R\cC^{(+)})_{AB}\big),
\end{equation}
and hence
\begin{equation}\label{eq:generalized-detector-inversion}
  \cC^{(+)}_{AB}
  =
  \frac{2}{\delta^2+\gamma^2}
  \big(\delta\,\Phi_{AB}-\gamma(R\Phi)_{AB}\big).
\end{equation}
Under the no-incoming detector identification \(\cC^{(+)}_{AB}=2\hat C_{AB}\),
\begin{equation}\label{eq:generalized-C-inversion}
  \hat C_{AB}
  =
  \frac{1}{\delta^2+\gamma^2}
  \big(\delta\,\Phi_{AB}-\gamma(R\Phi)_{AB}\big),
  \qquad
  \hat C_{AB}\hat C^{AB}
  =
  \frac{\Phi_{AB}\Phi^{AB}}{\delta^2+\gamma^2}.
\end{equation}
FAS eq. (VI.132) gives the same normalization invariantly:
\begin{equation}\label{eq:generalized-W-normalization}
  \frac{\Phi_{AB}\Phi^{AB}}{\delta^2+\gamma^2}
  =
  \frac14 W^+ .
\end{equation}
Thus the normalized tensor component
\begin{equation}\label{eq:Phi-normalized}
  \Phi^{\mathrm{norm}}_{AB}
  :=
  \frac{1}{\sqrt{\delta^2+\gamma^2}}\,\Phi_{AB}
\end{equation}
satisfies
\begin{equation}
  \Phi^{\mathrm{norm}}_{AB}\Phi^{\mathrm{norm}\,AB}
  =
  \frac14 W^+
  =
  \hat C_{AB}\hat C^{AB}.
\end{equation}
For the original umbilical transport one has \((\delta,\gamma)=(0,1)\), so \(\Phi^{\mathrm{norm}}=\Phi=\epsilon_B{}^E\hat C_{AE}\), recovering the previous construction.

If the divergence relation \(D^Bn^+_{AB}=-2N\epsilon_{AB}C^B\) and the no-incoming identification \(C^{(+)}_A=C_A\) remain valid, the vector detector is recovered exactly as before.  On a round cut define
\begin{equation}\label{eq:generalized-first-jet-news}
  \cN^{\mathrm{jet}}_{\delta,\gamma}
  :=
  \big(\Phi^{\mathrm{norm}}_{AB},\,\mathcal J C_A[n^+]\big) .
\end{equation}
Then
\begin{equation}
  \|\cN^{\mathrm{jet}}_{\delta,\gamma}\|^2
  =
  -\frac12 p_m,
  \qquad
  \cN^{\mathrm{jet}}_{\delta,\gamma}=0
  \iff
  n^+_{AB}=0
  \iff
  \text{no radiation at }\scri^+ .
\end{equation}
This extends the first-jet news from the single transport normalization \((\delta,\gamma)=(0,1)\) to the full nondegenerate generalized transport family.  It does not by itself remove the equipment, \(S^2\), no-incoming, or round-cut restrictions; those remain the hypotheses under which each algebraic step is valid.  The finite rotation-inversion algebra is formalized by \texttt{GeneralizedJetRotation.TensorRotation}, \texttt{sq\_inverse}, and \texttt{inverse\_zero\_iff}.

\subsection{Conformal weights and locality of the rank-two realization}

\statusD The obstruction to extending \(\mathcal J\) away from round cuts is already visible in conformal weights.  Under the residual rescaling \(q_{AB}\mapsto \tilde q_{AB}=\omega^2q_{AB}\) on a cut:

\begin{itemize}
  \item the FAS radiant news is gauge invariant as a covariant tensor, \(\tilde n^+_{AB}=n^+_{AB}\), and the lapse is invariant on \(\scri\), \(\tilde N=N\) by FAS eq. (E.8) at \(\Omega=0\);
  \item \(\tilde\epsilon_{AB}=\omega^2\epsilon_{AB}\), while \(\epsilon_A{}^B\) is invariant;
  \item for trace-free invariant \(n_{AB}\), the connection terms cancel in the divergence, so \(\widetilde{D^Bn_{AB}}=\omega^{-2}D^Bn_{AB}\);
  \item \(\tilde\Sigma_{AB}=\omega\Sigma_{AB}\) and \(\tilde{\bar m}^a=\omega^{-1}\bar m^a\), hence \(\dot n_{AB}-\Sigma_{(AC}n_{B)}{}^C\) has weight \(\omega^{-1}\).
\end{itemize}

The detector inversion
\begin{equation}\label{eq:detector-weight-div}
  D^Bn^+_{AB}=-2N\epsilon_{AB}C^B
\end{equation}
therefore forces
\begin{equation}\label{eq:detector-weights}
  C^B\mapsto \omega^{-4}C^B,
  \qquad
  C_A\mapsto \omega^{-2}C_A .
\end{equation}
On the other hand the FAS transport law and the invariant \(\epsilon_B{}^E\) give
\begin{equation}\label{eq:flux-weight}
  \Phi^+_{AB}\mapsto\omega^{-1}\Phi^+_{AB},
  \qquad
  \hat C_{AB}\mapsto\omega^{-1}\hat C_{AB}.
\end{equation}
Thus both detector squares have pointwise weight \(\omega^{-6}\):
\begin{equation}\label{eq:detector-square-weights}
  \hat C_{AB}\hat C^{AB}
  =
  q^{AC}q^{BD}\hat C_{AB}\hat C_{CD}
  \mapsto\omega^{-6}\hat C_{AB}\hat C^{AB},
  \qquad
  C_AC^A\mapsto\omega^{-6}C_AC^A,
\end{equation}
in agreement with \(p_m\mapsto\omega^{-6}p_m\).  The rank-two first-jet news can exist on arbitrary cuts only if there is a family of operators \(\mathcal J_q\) with the weighted covariance
\begin{equation}\label{eq:J-required-weight}
  \mathcal J_{\tilde q}(\omega^{-2}C)
  =
  \omega^{-1}\mathcal J_q(C),
\end{equation}
i.e. a map of conformal weight \(+1\) from weight-\(\omega^{-2}\) one-forms to weight-\(\omega^{-1}\) STF tensors.

To see why this is genuinely nonlocal, let \(\mathcal K_qC_{AB}=D_{(A}C_{B)}-\tfrac12q_{AB}D^DC_D\).  A direct conformal calculation gives, with \(\omega_A:=D_A\omega\),
\begin{equation}\label{eq:K-noncovariance}
  \mathcal K_{\tilde q}(\omega^{-2}C)_{AB}
  =
  \omega^{-2}\left[
    \mathcal K_q(C)_{AB}
    -2(\omega_AC_B+\omega_BC_A)
    +2q_{AB}\omega^DC_D
    \right].
\end{equation}
The bracket is trace-free but not equal to \(\mathcal K_q(C)\); the anomaly is first order in \(\omega\).  More generally, any natural local tensor differential operator in two dimensions is built by contracting \(q_{AB}\), \(q^{AB}\), \(\epsilon_{AB}\), \(\epsilon^{AB}\), curvature tensors, and covariant derivatives.  Each contraction with \(q\), \(\epsilon\), or curvature shifts the conformal weight by an even integer, while covariant derivatives do not change the shift.  Hence every such operator shifts conformal weight by an element of \(2\mathbb Z\).  The required shift is
\begin{equation}
  (-1)-(-2)=1\notin2\mathbb Z .
\end{equation}
Therefore no local tensor differential operator, of any finite order, can realize the required \(\mathcal J_q\).  This rules out a curvature-corrected local version of \(\mathcal K_q\), but it does not exclude a pseudodifferential operator: the round-cut \(\mathcal J=(2/(-\Delta-2))^{1/2}\mathcal K\) has order \(-1\) and exactly the required weight shift.  The precise geometric problem is a conformally covariant operator \(\mathcal J_q\) satisfying \eqref{eq:J-required-weight}, annihilating the conformal-Killing kernel, and preserving
\begin{equation}\label{eq:J-covariant-isometry}
  \int_{\tilde S}\mathcal J_{\tilde q}(\omega^{-2}C)\,\mathcal J_{\tilde q}(\omega^{-2}C)\,d\tilde S
  =
  \omega^{-4}
  \int_S\mathcal J_q(C)\,\mathcal J_q(C)\,dS .
\end{equation}
The finite weight algebra and the parity obstruction are formalized by \texttt{ConformalWeightObstruction.LocalTensorOperator}, \texttt{tensor\_square\_weight}, \texttt{vector\_square\_weight}, and \texttt{no\_local\_rank\_two\_map}.

\subsection{Metric-dependent arbitrary-cut isometry}

\statusD The round-cut restriction in the vector-to-tensor step can be removed if one asks for metric dependence rather than conformal covariance.  Let \((S,q)\) be any closed oriented Riemannian two-surface.  Define the conformal Killing operator
\begin{equation}\label{eq:arbitrary-cut-K}
  \mathcal K_q v_{AB}
  :=
  D_{(A}v_{B)}-\frac12q_{AB}D^Cv_C ,
  \qquad
  v_A\in\Omega^1(S).
\end{equation}
Its kernel is exactly the space of conformal Killing one-forms.  On a round \(S^2\) this is the harmonic \(l=1\) space, with the constant mode absent for one-forms.  The operator is elliptic, so
\begin{equation}
  P_q:=\mathcal K_q^\ast\mathcal K_q
\end{equation}
is a nonnegative self-adjoint elliptic operator with finite-dimensional kernel \(\ker\mathcal K_q\).  Let \(P_q^{-1/2}\) denote its inverse square root on \((\ker\mathcal K_q)^\perp\), extended by zero on \(\ker\mathcal K_q\), and define
\begin{equation}\label{eq:arbitrary-cut-J}
  \widehat{\mathcal J}_q
  :=
  \mathcal K_q P_q^{-1/2}.
\end{equation}
For every \(v\in(\ker\mathcal K_q)^\perp\),
\begin{equation}\label{eq:arbitrary-cut-isometry}
  \int_S
  \widehat{\mathcal J}_q v_{AB}\,
  \widehat{\mathcal J}_q v^{AB}\,dS
  =
  \int_S v_Av^A\,dS ,
\end{equation}
and \(\widehat{\mathcal J}_qv=0\iff v=0\).  On all one-forms, \(\ker\widehat{\mathcal J}_q=\ker\mathcal K_q\).  On the unit round sphere the eigenvalues of \(P_q\) on the \(l\ge2\) vector harmonics are \((L-2)/2\), so \(\widehat{\mathcal J}_q=\mathcal K_q\sqrt{2/(-\Delta-2)}\), recovering the operator in \eqref{eq:rank2-pair-news}.  Thus \eqref{eq:arbitrary-cut-J} is the canonical polar-decomposition extension of the round-cut Hodge inverse to arbitrary cuts.

Assume now that the detector one-form recovered from the outgoing news satisfies
\begin{equation}\label{eq:CKV-orthogonality}
  C_A[n^+]\in(\ker\mathcal K_q)^\perp .
\end{equation}
Then the first-jet news can be defined on any closed oriented cut by
\begin{equation}\label{eq:arbitrary-cut-first-jet}
  \cN^{\mathrm{jet}}_{q,\delta,\gamma}
  :=
  \big(\Phi^{\mathrm{norm}}_{AB},\,\widehat{\mathcal J}_q C_A[n^+]\big)
  \in
  \mathfrak{STF}(S)\oplus\mathfrak{STF}(S),
\end{equation}
with generalized \((\delta,\gamma)\) normalization as in \eqref{eq:generalized-first-jet-news}.  The isometry identity \eqref{eq:arbitrary-cut-isometry} gives
\begin{equation}
  \|\cN^{\mathrm{jet}}_{q,\delta,\gamma}\|_{L^2}^2
  =
  -\frac12\int_S p_m\,dS ,
  \qquad
  \cN^{\mathrm{jet}}_{q,\delta,\gamma}=0
  \iff
  \hat C_{AB}=0,\ C_A=0 .
\end{equation}
This removes the round-sphere restriction from the \emph{norm and zero-set} realization of the first-jet news.  The intrinsic replacement for the discarded \(l=0,1\) modes is the conformal-Killing kernel, which is now handled by orthogonal projection rather than by round harmonics.

The operator \(\widehat{\mathcal J}_q\) above is built from the unweighted \(L^2\) adjoint, so its construction is metric dependent and does not by itself display conformal covariance.  The following weighted refinement supplies the missing intertwiner.

\subsection{Conformally covariant weighted vector-to-tensor intertwiner}

\statusD Work on an oriented cut with \(S^2\) topology.  Let
\(\mu_q:=(\det q)^{1/4}\) be the natural positive one-density, so that under
\(q\mapsto\tilde q=\omega^2q\) one has \(\tilde\mu=\omega\mu\).  Say that a covector has conformal weight \(w\) when \(\tilde v_A=\omega^w v_A\).  The detector \(C_A\) has weight \(w=-2\), and the target STF tensor has weight \(w=-1\).

Consider the conformal Killing operator on weight-\(2\) one-forms,
\begin{equation}\label{eq:weighted-K}
  \mathcal K_q^{(2)}v_{AB}
  :=
  D_{(A}v_{B)}-\frac12q_{AB}D^Cv_C ,
  \qquad v_A\in\mathcal E_1[2].
\end{equation}
A direct Weyl-variation computation gives
\begin{equation}\label{eq:weighted-K-covariance}
  \mathcal K_{\tilde q}^{(2)}(\omega^2v)
  =
  \omega^2\mathcal K_q^{(2)}v .
\end{equation}
This is the distinguished weight at which \(\mathcal K_q\) is conformally covariant; the unweighted operator in \eqref{eq:arbitrary-cut-K} fails to be covariant.  To take an adjoint with the correct weight, use the pairings
\begin{align}\label{eq:weighted-pairings}
  \langle v,u\rangle_{V,q}
  &:=
  \int_S q^{AB}v_Au_B\,\mu_q^{-8}\,dS_q ,
  &
  \langle S,U\rangle_{W,q}
  &:=
  \int_S q^{AC}q^{BD}S_{AB}U_{CD}\,\mu_q^{-6}\,dS_q .
\end{align}
If \(v=\mu_q^4C\) and \(S=\mu_q^3T\), these are exactly the geometric norms
\(\int|C|^2\,dS_q\) and \(\int|T|^2\,dS_q\).  Under
\(q\mapsto\tilde q=\omega^2q\), with \(v\mapsto\omega^2v\) and
\(S\mapsto\omega^2S\), both pairings acquire the same factor \(\omega^{-4}\).  Hence the formal adjoint of \(\mathcal K_q^{(2)}\) with respect to \eqref{eq:weighted-pairings} obeys
\begin{equation}\label{eq:weighted-K-adjoint-covariance}
  \mathcal K_{\tilde q}^{(2)\dagger}(\omega^2S)
  =
  \omega^2\mathcal K_q^{(2)\dagger}S .
\end{equation}
Define the elliptic nonnegative operator on weight-\(2\) one-forms
\begin{equation}\label{eq:weighted-P}
  P_q^{(2)}
  :=
  \mathcal K_q^{(2)\dagger}\mathcal K_q^{(2)} ,
\end{equation}
with adjoint taken in \eqref{eq:weighted-pairings}.  Combining
\eqref{eq:weighted-K-covariance} and
\eqref{eq:weighted-K-adjoint-covariance} gives
\begin{equation}\label{eq:weighted-P-covariance}
  P_{\tilde q}^{(2)}(\omega^2v)
  =
  \omega^2P_q^{(2)}v .
\end{equation}
Its kernel is the conformal-Killing kernel of \(\mathcal K_q^{(2)}\), and its inverse square root on the orthogonal complement inherits the same covariance,
\begin{equation}\label{eq:weighted-P-inverse-covariance}
  (P_{\tilde q}^{(2)})^{-1/2}(\omega^2v)
  =
  \omega^2(P_q^{(2)})^{-1/2}v .
\end{equation}
Therefore define the conformally covariant weighted intertwiner
\begin{equation}\label{eq:covariant-J}
  \mathcal J_q^{\mathrm{cov}}C
  :=
  \mu_q^{-3}\,
  \mathcal K_q^{(2)}
  (P_q^{(2)})^{-1/2}
  \big(\mu_q^4 C\big) .
\end{equation}
The weights are \(C\in\mathcal E_1[-2]\),
\(\mu_q^4C\in\mathcal E_1[2]\),
\((P_q^{(2)})^{-1/2}(\mu_q^4C)\in\mathcal E_1[2]\),
\(\mathcal K_q^{(2)}(\cdots)\in\mathcal E_{(AB)_0}[2]\), and multiplication by
\(\mu_q^{-3}\) lands in \(\mathcal E_{(AB)_0}[-1]\).  Using the covariance laws and \(\tilde C=\omega^{-2}C\),
\begin{align}
  \mathcal J_{\tilde q}^{\mathrm{cov}}(\omega^{-2}C)
  &=
  \tilde\mu^{-3}
  \mathcal K_{\tilde q}^{(2)}
  (P_{\tilde q}^{(2)})^{-1/2}
  \big(\tilde\mu^4\omega^{-2}C\big)
  \notag\\
  &=
  \omega^{-3}\mu^{-3}
  \mathcal K_{\tilde q}^{(2)}
  (P_{\tilde q}^{(2)})^{-1/2}
  \big(\omega^2\mu^4C\big)
  \\
  &=
  \omega^{-3}\mu^{-3}
  \mathcal K_{\tilde q}^{(2)}
  \big(\omega^2(P_q^{(2)})^{-1/2}(\mu^4C)\big)
  =
  \omega^{-1}\mathcal J_q^{\mathrm{cov}}C ,
  \notag
\end{align}
which is exactly the required law \eqref{eq:J-required-weight}.

Because \(P_q^{(2)}=\mathcal K_q^{(2)\dagger}\mathcal K_q^{(2)}\), the operator
\(\mathcal K_q^{(2)}(P_q^{(2)})^{-1/2}\) is an isometry from the \(V\)-pairing to the \(W\)-pairing on \((\ker\mathcal K_q^{(2)})^\perp\).  Unwinding the definitions in \eqref{eq:weighted-pairings}, this says that for every detector \(C\) whose weighted representative \(\mu_q^4C\) has no conformal-Killing component,
\begin{equation}\label{eq:covariant-J-isometry}
  \int_S
  \mathcal J_q^{\mathrm{cov}}C_{AB}\,
  \mathcal J_q^{\mathrm{cov}}C^{AB}\,dS_q
  =
  \int_S C_AC^A\,dS_q ,
  \qquad
  \mathcal J_q^{\mathrm{cov}}C=0\iff C=0 .
\end{equation}
No uniformization or round-cut specialization is needed for this identity; the conformal weight bookkeeping alone transfers the polar-decomposition isometry to arbitrary cuts.

\subsubsection{Canonical polar uniqueness of the intertwiner}

\statusD The operator \eqref{eq:covariant-J} is not an arbitrary choice.  Define the local weighted conformal Killing operator between the geometric Hilbert spaces of detector one-forms and STF tensors,
\begin{equation}\label{eq:weighted-local-K}
  \widetilde{\mathcal K}_q
  :=
  \mu_q^{-3}\,\mathcal K_q^{(2)}\,\mu_q^4 :
  \mathcal E_1[-2]\longrightarrow \mathcal E_{(AB)_0}[-1] .
\end{equation}
By \eqref{eq:weighted-K-covariance} and the density weights,
\begin{equation}\label{eq:weighted-local-K-covariance}
  \widetilde{\mathcal K}_{\tilde q}(\omega^{-2}C)
  =
  \omega^{-1}\widetilde{\mathcal K}_qC .
\end{equation}
Thus the obstruction to a local map is not a defect of \eqref{eq:weighted-local-K}; it is the fact that the desired news must be an \emph{isometric} realization, while \(\widetilde{\mathcal K}_q\) itself is not isometrically norm preserving.  Let \(\widetilde{\mathcal K}_q=\mathcal U_q|\widetilde{\mathcal K}_q|\) be its polar decomposition on
\(V_q:=(\ker\widetilde{\mathcal K}_q)^\perp\).  Since both geometric norms scale by the same factor \(\omega^{-4}\), the rescaled maps are unitary after the common normalization by \(\omega^2\).  Uniqueness of the positive square root and of the polar factor therefore gives
\begin{align}
  |\widetilde{\mathcal K}_{\tilde q}|(\omega^{-2}C)
  &:=
  \omega^{-2}|\widetilde{\mathcal K}_q|C ,\\
  \mathcal U_{\tilde q}(\omega^{-2}C)
  &:=
  \omega^{-1}\mathcal U_qC .
\end{align}
Moreover, unwinding the definitions of the adjoint in \eqref{eq:weighted-pairings} identifies
\begin{equation}\label{eq:polar-factor-J}
  \mathcal U_q
  =
  \widetilde{\mathcal K}_q|\widetilde{\mathcal K}_q|^{-1}
  =
  \mu_q^{-3}\mathcal K_q^{(2)}
  (P_q^{(2)})^{-1/2}\mu_q^4
  =
  \mathcal J_q^{\mathrm{cov}} .
\end{equation}
Consequently \(\mathcal J_q^{\mathrm{cov}}\) is the \emph{unique} isometric conformally covariant map on the non-kernel sector factorization of \(\widetilde{\mathcal K}_q\): if \(\mathcal U'_q:V_q\to W_q\) is another isometry with \(\widetilde{\mathcal K}_q=\mathcal U'_q|\widetilde{\mathcal K}_q|\), then \(\mathcal U'_q=\mathcal J_q^{\mathrm{cov}}\).  This turns the weight-\(+1\) construction from an existence statement into a canonical choice.  The finite algebraic core is formalized by \texttt{ConformalJetIntertwiner.PolarAlgebra} and \texttt{polar\_unique}.

Thus the arbitrary-cut first-jet news can be defined by
\begin{equation}\label{eq:covariant-first-jet}
  \cN^{\mathrm{jet},\mathrm{cov}}_{q,\delta,\gamma}
  :=
  \big(
    \Phi^{\mathrm{norm}}_{AB},
    \mathcal J_q^{\mathrm{cov}}C_A[n^+]
  \big) ,
\end{equation}
and it obeys
\[
  \|\cN^{\mathrm{jet},\mathrm{cov}}_{q,\delta,\gamma}\|_{L^2}^2
  =
  -\frac12\int_S p_m\,dS_q ,
  \qquad
  \cN^{\mathrm{jet},\mathrm{cov}}_{q,\delta,\gamma}=0
  \iff
  \hat C_{AB}=0,\ C_A=0 ,
\]
provided the recovered detector one-form has no conformal-Killing component.  For the inverse-divergence source from an STF news tensor on \(S^2\), this is exactly the standard Codazzi solvability condition.  Hence the round-cut restriction and the conformal-covariance obstruction are both removed in the controlled outgoing sector.  The remaining gap is no longer cut covariance but the coupling to AldS/symplectic flux, extension beyond strong outgoing equipment, and the unresolved flat-limit coefficient \(\beta_+\).  The finite weight algebra is formalized by \texttt{ConformalJetIntertwiner.Algebra}, \texttt{input\_weight}, \texttt{output\_weight}, and \texttt{jet\_square\_scaling}.

\subsection{Explicit Codazzi inversion on a round cut}

\statusD It is useful to see exactly how nonlocal the proposed news must be.  Work in the conformal gauge in which the cut metric \(q_{AB}\) is the unit round metric, and isolate the simplified Codazzi equation
\begin{equation}\label{eq:codazzi-simplified}
  D_{[C}M_{A]B}=2q_{B[C}X_{A]}
  \quad\Longleftrightarrow\quad
  D^A M_{AB}=2X_B ,
\end{equation}
with \(M_{AB}=M_{(AB)}\), \(q^{AB}M_{AB}=0\).  For a scalar harmonic \(Y_{lm}\) with \(\Delta Y_{lm}=-l(l+1)Y_{lm}\), define the STF potentials
\begin{equation}
  H_{AB}(Y)=D_AD_BY-\frac12q_{AB}\Delta Y,
  \qquad
  G_{AB}(Y)=\epsilon_{(A}{}^{C}D_{B)}D_CY .
\end{equation}
Commuting covariant derivatives on the round sphere gives
\begin{align}
  D^A H_{AB}(Y_{lm})
  &=
  \left(1-\frac{l(l+1)}2\right)D_BY_{lm},\\
  D^A G_{AB}(Y_{lm})
  &=
  \frac{1-l(l+1)}2\,\epsilon_B{}^C D_CY_{lm}.
\end{align}
Write the Hodge potentials of the source as
\begin{equation}
  X_A=D_AX+\epsilon_A{}^B D_Bx,
  \qquad
  X=\sum_{l\ge2,m}X_{lm}Y_{lm},
  \quad
  x=\sum_{l\ge2,m}x_{lm}Y_{lm},
\end{equation}
where the \(l=1\) terms are absent by \eqref{eq:CKV-integrability} and the \(l=0\) terms are gauge.  Then the unique trace-free solution of \eqref{eq:codazzi-simplified} is
\begin{equation}\label{eq:codazzi-inversion}
  M_{AB}
  =
  \sum_{l\ge2,m}
  \left[
    \frac{2}{1-l(l+1)/2}\,X_{lm}\,H_{AB}(Y_{lm})
    +
    \frac{4}{1-l(l+1)}\,x_{lm}\,G_{AB}(Y_{lm})
  \right].
\end{equation}
Both denominators are nonzero for \(l\ge2\), and both potential operators vanish at \(l=0,1\); this is the harmonic form of the uniqueness of trace-free Codazzi tensors on \(S^2\).

Equation \eqref{eq:codazzi-inversion} is an inverse elliptic operation: the candidate \(M_{AB}\) is obtained from \(X_A\) by applying harmonic-space multipliers of order \(l^{-2}\), followed by two derivatives.  Thus the cut-news candidate is intrinsically nonlocal on the cut, even in the round gauge.  The inversion coefficients contain no \(\Lambda\); in this gauge any singular \(\Lambda\to0\) behavior must enter through the source \(X_A\), the conformal representative, or the extrinsic-curvature terms, not through the Codazzi operator itself.  The full source displayed in \eqref{eq:X-source-div} differs from the simplified right-hand side of \eqref{eq:codazzi-simplified} by the Cotton and extrinsic-curvature terms shown there; the same inversion applies after those known terms are absorbed into \(X_A\).

For the suggested ansatz \(X_A=\Delta f\,D_Af\), the potentials obey
\begin{equation}
  \Delta X=(\Delta f)^2+D^A(\Delta f)\,D_Af,
  \qquad
  \epsilon^{AB}D_AX_B=-\Delta x
  =-\epsilon^{AB}D_A(\Delta f)\,D_Bf ,
\end{equation}
so \eqref{eq:codazzi-inversion} gives an explicit formal \(M_{AB}[f]\) in the round gauge.  Because \(X_A=0\) iff \(f\) is constant, this formal candidate has the source-zero set required by the program; it still has no established relation to the AldS charge flux or to the outgoing positivity density \eqref{eq:outgoing-pm}.

\subsection{Codazzi uniqueness and the zero-set condition}

On \(S^2\), trace-free Codazzi tensors vanish identically.  Therefore, once a solution of \(D_{[C}N_{A]B}=X_{CAB}\) exists,
\begin{equation}
  X_A=0 \text{ on } S
  \iff
  D_{[C}N_{A]B}=0
  \iff
  N_{AB}=0 .
\end{equation}
Thus the paper's target condition \(X_A=0\iff N_{AB}=0\) is automatic after solvability.  The genuinely open requirement is to choose \(X_A\) so that its zero set matches the physical no-radiation detector \([C,\cF]=0\).

\section{Obstructions to pointwise linear news squares}

\subsection{Scaling of the detector}

A tensor built only from the STF part of the super-Poynting vector cannot be the AldS news.  The detector \(p_i\) is linear in the electric Weyl/stress tensor \(D_{ij}\); a tensor obtained linearly from \(p_i\) is at least linear in \(D\), and its square is at least quadratic in \(D\).  The AldS flux \eqref{eq:aldS-flux}, however, is linear in \(\cT\), hence linear in \(D\).  This rules out the most naive “dual of the detector” ans{\"a}tze.

\subsection{A finite-dimensional no-go theorem}

Under \(Q_-^\mu=0\), the positive density is
\begin{equation}
  -p_m=2(\hat C_{AB}\hat C^{AB}+C_A C^A)\ge0,
\end{equation}
where \(\hat C_{AB}\) and \(C_A\) are independent.  Consider a pointwise linear ansatz
\begin{equation}
  N_{AB}=L(\hat C)_{AB}+K(C)_{AB}
\end{equation}
into the two-dimensional space of STF tensors on the cut.

\begin{theorem}[\statusD]
Let \(E\) be the two-dimensional real vector space of STF cut tensors with Euclidean norm induced by \(q_{AB}\).  There do not exist linear maps \(L,K:E\to E\) such that
\begin{equation}
  \|Lx\|=\|x\|,\qquad \|Ky\|=\|y\|,
\end{equation}
and
\begin{equation}
  \|Lx+Ky\|^2=\|x\|^2+\|y\|^2
\end{equation}
for all \(x,y\in E\).
\end{theorem}
\begin{proof}
The first condition makes \(L\) injective; finite-dimensionality makes it surjective.  Choose \(y\neq0\).  Pick \(x\) with \(Lx=Ky\).  Then
\begin{equation}
  \|Lx+Ky\|^2=\|2Ky\|^2=4\|Ky\|^2,
\end{equation}
whereas
\begin{equation}
  \|x\|^2+\|y\|^2=\|Ky\|^2+\|Ky\|^2=2\|Ky\|^2,
\end{equation}
a contradiction since \(Ky\neq0\).
\end{proof}

The Lean file proves this as \texttt{NewsSquareObstruction.no\_linear\_pointwise\_square}.  The theorem rules out a purely pointwise linear square of the outgoing Cotton data.  It does not by itself rule out the cut-Codazzi construction \(N_{AB}=V_{AB}+M_{AB}\); the following result sharpens the nonlocal case for maps from the full detector pair into one rank-two tensor.

\subsection{A representation-theoretic nonlocal obstruction}

\statusD The pointwise dimension count extends to a nonlocal obstruction on the round cut.  Let \(\mathcal H_l\) denote the \((2l+1)\)-dimensional irreducible \(\mathrm{SO}(3)\)-representation of scalar harmonics of degree \(l\ge2\).  On the unit round sphere,
\begin{equation}
  \mathfrak{STF}(S^2)_l
  =
  \operatorname{span}\{H_{AB}(Y_{lm}),G_{AB}(Y_{lm})\}_{m}
  \cong
  \mathcal H_l\oplus\mathcal H_l ,
\end{equation}
while vector harmonics decompose as
\begin{equation}
  \mathfrak{X}(S^2)_l
  =
  \operatorname{span}\{D_AY_{lm},\epsilon_A{}^BD_BY_{lm}\}_{m}
  \cong
  \mathcal H_l\oplus\mathcal H_l .
\end{equation}
Consequently the outgoing detector pair \((\hat C_{AB},C_A)\) lies in \(4\mathcal H_l\), whereas a single rank-two STF news tensor lies in \(2\mathcal H_l\).

\begin{proposition}[\statusD]
Let \(L:\mathfrak{STF}(S^2)\oplus\mathfrak{X}(S^2)\to\mathfrak{STF}(S^2)\) be a linear map that is \(\mathrm{SO}(3)\)-equivariant, possibly nonlocal and with arbitrary spectral multipliers.  Then \(L\) has a nontrivial kernel on every \(l\ge2\) sector.  In particular, no such \(L\) can satisfy
\begin{equation}
  \|L(x,y)\|^2=\|x\|^2+\|y\|^2
\end{equation}
for all detector pairs \((x,y)\).
\end{proposition}
\begin{proof}
Schur's lemma reduces the restriction of \(L\) to each isotypic component to a matrix
\begin{equation}
  L_l:\mathbb R^4\otimes\mathcal H_l\longrightarrow\mathbb R^2\otimes\mathcal H_l .
\end{equation}
The domain has fibre dimension \(4\) and the target fibre dimension \(2\), so \(L_l\) cannot be injective.  Norm preservation would imply injectivity, a contradiction.
\end{proof}

Thus the full detector pair cannot be compressed into one rotation-covariant rank-two news tensor by any linear nonlocal operation on a round cut.  The minimal linear target preserving \(-p_m\) is the pair itself.  A rank-two object such as \(n^+_{AB}=2V_{AB}\) survives only because the strongly equipped outgoing sector restricts attention to a two-dimensional subdetector, namely the Cotton one-form that generates \(V_{AB}\) by inverse divergence.  The obstruction assumes linearity and \(\mathrm{SO}(3)\)-equivariance; nonlinear constructions, symmetry-breaking choices, \(\xi\)-dependent targets, and multi-cut or time-dependent maps are not excluded.

The finite-mode algebra is \texttt{HarmonicDimensionObstruction.no\_injective\_pair\_to\_tensor}: no injective \(\mathbb R\)-linear map exists from a \(4\)-dimensional mode to a \(2\)-dimensional mode.

\section{Linearized de Sitter waves}

\subsection{Variation of the boundary Cotton and stress tensors}

\statusD Work on the round \(S^3\) boundary of de Sitter with
\begin{equation}
  h_{ij}=\ell^2(q_{ij}+\varepsilon\kappa_{ij}),
  \qquad
  D^i\kappa_{ij}=0,
  \qquad
  q^{ij}\kappa_{ij}=0 .
\end{equation}
Using the FG relation \(D_{ij}\propto h^{(3)}_{ij}\) and \(C_{ij}\propto -\ell^2Y_{ij}\), a direct variation on the unit round \(S^3\) gives
\begin{equation}
  \delta P_{ij}
  =
  -\frac12(\Delta+1)\kappa_{ij},
  \qquad
  \delta Y_{ij}
  =
  -\frac12
  \epsilon_{(i}{}^{kl}
  D_k(\Delta+1)\kappa_{j)l}.
\end{equation}
The detector is bilinear:
\begin{equation}\label{eq:detector-bilinear}
  \bar{\cP}_i
  \propto
  3\ell^2
  \epsilon_i{}^{jk}
  Y_{j\ell}h^{(3)\ell}{}_k .
\end{equation}

Consequences:
\begin{itemize}
  \item Around pure de Sitter, \(Y^{(0)}=D^{(0)}=0\), so \(Y,D=O(\varepsilon)\) and \(\bar{\cP}=O(\varepsilon^2)\).  The super-Poynting detector has no first-order part and cannot itself be linear Bondi news.
  \item Around a nonradiating Coulomb background such as Kottler, \(Y^{(0)}=0\) but \(D^{(0)}\neq0\), so
  \begin{equation}
    \bar{\cP}^{(1)}_i
    \propto
    3\ell^2
    \epsilon_i{}^{jk}
    \delta Y_{j\ell}D^{(0)\ell}{}_k
  \end{equation}
  can be linear in the wave amplitude.
\end{itemize}
The bilinear algebra is recorded in Lean as \texttt{commutator\_scaleSym\_scaleSym}, \texttt{superPoynting\_scaleSym\_scaleSym}, and \texttt{linearized\_superPoynting\_is\_quadratic}.

\subsection{Explicit TT modes and Hamiltonian fluxes}

\statusL Ashtekar--Bonga--Kesavan solve the linearized field equations on the Poincaré patch of de Sitter \cite{ashtekar2015linear}.  Writing \(g_{ab}=a^2(\eta)(\mathring g_{ab}+\epsilon h_{ab})\), \(a=-1/(H\eta)\), \(\Omega=-H\eta\), TT modes satisfy
\begin{equation}
  \mathring{\Box}h_{ab}+\frac{2}{\eta}h'_{ab}=0,
  \qquad
  h_{\vec k}^{(s)\prime\prime}-\frac{2}{\eta}h_{\vec k}^{(s)\prime}+k^2h_{\vec k}^{(s)}=0,
\end{equation}
with solution
\begin{equation}
  h_{\vec k}^{(s)}
  =
  -2H\left[
    E_{\vec k}^{(s)}
    \left(\eta\cos k\eta-k^{-1}\sin k\eta\right)
    -
    B_{\vec k}^{(s)}
    \left(\eta\sin k\eta+k^{-1}\cos k\eta\right)
  \right].
\end{equation}
The rescaled Weyl data at \(\scri^+\) are
\begin{equation}
  \cE_{ab}|_{\scri^+}
  =
  k^2E_{\vec k}^{(s)}\mathfrak e^{(s)}_{ab},
  \qquad
  \cB_{ab}|_{\scri^+}
  =
  k^2B_{\vec k}^{(s)}{}^\ast\mathfrak e^{(s)}_{ab}.
\end{equation}
Gauge-invariant quadratic charges include
\begin{equation}
  \cH_T
  =
  \frac{1}{2H\kappa}
  \int_{\scri^+}
  \cE_{cd}
  \left(
    \cL_T h_{ab}+2Hh_{ab}
  \right)
  \mathring q^{ac}\mathring q^{bd}.
\end{equation}
The \(B_{\vec k}^{(s)}=0\) branch preserves conformal flatness of the first-order boundary geometry and carries zero energy, momentum, and angular momentum.  These quadratic Hamiltonian fluxes are robust linear-theory objects, not nonlinear news tensors.

In Fourier modes the energy flux takes the form \cite{ashtekar2015linear}
\begin{equation}\label{eq:linear-energy-fourier}
  \cH_T
  =
  \frac{H}{\kappa}
  \int\frac{\diff^3k}{(2\pi)^3}\,k
  \sum_{s=1}^{2}
  \left(
    E_{\vec k}^{(s)}\mathcal L_k(B_{\vec k}^{(s)})^\star
    +
    2E_{\vec k}^{(s)}(B_{\vec k}^{(s)})^\star
  \right).
\end{equation}
The electric and magnetic boundary amplitudes \(E_{\vec k}^{(s)}\) and \(B_{\vec k}^{(s)}\) are independent coordinates on the true linear phase space.  Equation \eqref{eq:linear-energy-fourier} is linear in \(E\) for fixed \(B\).  Hence the admissible change \(E\mapsto-E\), \(B\mapsto B\) reverses the sign of \(\cH_T\).

\begin{proposition}[\statusD: sign-flip obstruction to a linearized news square]
There is no nonnegative density \(\sigma\) on the full linearized TT phase space such that \(\cH_T=\pm\int\sigma^2\diff\mu\) for every solution.  More precisely, any such representation would force \(\cH_T=0\) on the entire phase space.
\end{proposition}
\begin{proof}
A square integral is sign-definite.  If the proposed representation holds for a solution with \(\cH_T\neq0\), applying it to the solution with \(E\mapsto-E\) gives the opposite sign for the same nonnegative right-hand side, a contradiction.  Thus the only sign-definite representation is the trivial one.
\end{proof}

This does not say that outgoing, source-generated linear waves have negative flux.  It says the restriction to such waves is essential: only after a no-incoming condition relates \(B\) nonlocally to \(E\) can a positive quadratic energy become possible.  Since \eqref{eq:linear-energy-fourier} is linear in the generator and bilinear in phase-space data, the natural linearized object is a symplectic pairing, not a single local square.  The finite algebraic core is formalized in Lean as \texttt{LinearizedFluxObstruction.SignFlipNewsModel}, \texttt{energy\_zero\_of\_sign\_flip\_square}, and \texttt{no\_sign\_flip\_square}.

\subsection{Finite de Sitter charges do not detect all news}

\statusL The ABK flux formulas show a structural feature beyond sign indefiniteness.  At \(\scri^+\), the electric boundary Weyl field \(\cE_{ab}\) is determined by the Fourier coefficients \(E_{\vec k}^{(s)}\), while the surviving boundary metric perturbation \(h_{ab}\) is determined by \(B_{\vec k}^{(s)}\).  Every Hamiltonian flux for a Poincaré generator is a cross term:
\begin{equation}\label{eq:ABK-cross-flux}
  \cH_S
  =
  \frac{1}{2H\kappa}
  \int_{\scri^+}
  \cE_{ab}\,\mathpounds_S h_{cd}\,\mathring q^{ac}\mathring q^{bd},
  \qquad
  \cH_T
  =
  \frac{1}{2H\kappa}
  \int_{\scri^+}
  \cE_{cd}\big(\mathpounds_T h_{ab}+2Hh_{ab}\big)
  \mathring q^{ac}\mathring q^{bd}.
\end{equation}
Thus the six momentum/rotation charges and the energy are all bilinear in \(E\) and \(B\).  ABK explicitly note that the energy has an infinite-dimensional zero-energy subspace \cite{ashtekar2015linear}.

\statusD A sharper statement follows by Fourier support separation.  Choose \(E_{\vec k}^{(s)}\) and \(B_{\vec k}^{(s)}\) smooth, nonzero, and compactly supported in disjoint open subsets of momentum space.  Since \(\mathpounds_S\) and \(\mathpounds_T\) act by differential/dilation operators in \(\vec k\), they preserve momentum support.  The products in \eqref{eq:ABK-cross-flux} then integrate to zero for every generator, so
\begin{equation}\label{eq:all-Poincare-flux-zero}
  \cH_K=0
  \quad\text{for all seven Poincaré generators }K .
\end{equation}
At the same time \(\cE_{ab}\neq0\) and \(\cB_{ab}\neq0\).  For generic polarizations their position-space commutator is nonzero, so the detector pair, and hence the first-jet news \eqref{eq:first-jet-news}, is nonzero:
\begin{equation}
  \cH_K=0\ \forall K
  \quad\nRightarrow\quad
  \cN^{\mathrm{jet}}=0 .
\end{equation}
The opposite implication does hold in the overlap of the linearized no-incoming sector and the strongly equipped outgoing sector: \(\cN^{\mathrm{jet}}=0\) gives \(\hat C_{AB}=C_A=0\), hence the radiative projection of the linearized Hamiltonian news pair vanishes and \(\Pi_{\mathrm{rad}}\mathcal F_\xi=0\) for every allowed generator.  Coulombian charges may remain nonzero but are conserved.

This is a charge-blindness theorem, not a contradiction with the unified zero set: detector-zero is a strong no-radiation statement, whereas vanishing of the finite family of de Sitter charge fluxes is much weaker.  It also explains why the central theorem cannot be phrased as equality between a news square and any individual AldS charge flux.  A viable flux representation must be fixed-\(\xi\), \(\xi\)-weighted, or bilinear, and must use news data beyond the seven finite charge evaluations.  The finite logical core is formalized by \texttt{ChargeFluxBlindness.Model}, \texttt{news\_zero\_flux\_zero}, and \texttt{not\_flux\_zero\_imp\_news\_zero}.

\subsection{Conditional symplectic dual-news representation}

\statusD The charge-blindness result suggests the correct flux law is not a news square but a pairing with a generator-dependent dual news.  This can be made precise in the controlled linear sector.  Let \(\Gamma_{\mathrm{rad}}\) denote the no-incoming outgoing radiative phase space, let \(\Omega\) be its symplectic form, and let
\begin{equation}
  \cN^{\mathrm{jet}}_{\mathrm{rad}}\colon \Gamma_{\mathrm{rad}}\longrightarrow \cH_S
  :=
  \mathfrak{STF}(S)\oplus\mathfrak{STF}(S)
\end{equation}
be the first-jet news map into the Hilbert space of rank-two pairs.  Assume:

\begin{enumerate}
  \item \(\cN^{\mathrm{jet}}_{\mathrm{rad}}\) is linear, injective and of closed range on \(\Gamma_{\mathrm{rad}}\);
  \item for every allowed generator \(\xi\), the Hamiltonian flow \(X_\xi\) preserves \(\Gamma_{\mathrm{rad}}\), and the radiative flux one-form
    \begin{equation}
      \mathcal F_\xi^{\mathrm{rad}}(\delta)
      :=
      \Omega(\delta,X_\xi),
      \qquad \delta\in\Gamma_{\mathrm{rad}},
    \end{equation}
    is bounded in the \(\cH_S\)-norm;
  \item \(\cN^{\mathrm{jet}}_{\mathrm{rad}}\) intertwines the residual gauge action with its natural action on \(\cH_S\).
\end{enumerate}

Define a linear functional on the range by
\begin{equation}
  \ell_\xi\big(\cN^{\mathrm{jet}}_{\mathrm{rad}}(\delta)\big)
  :=
  \mathcal F_\xi^{\mathrm{rad}}(\delta) .
\end{equation}
Injectivity makes \(\ell_\xi\) well defined, and boundedness plus the Riesz representation theorem give a unique dual news field \(\mathscr N_\xi\in\overline{\operatorname{Ran}\cN^{\mathrm{jet}}_{\mathrm{rad}}}\) such that
\begin{equation}\label{eq:dual-news-flux}
  \boxed{
    \mathcal F_\xi^{\mathrm{rad}}(\delta)
    =
    \big\langle
      \cN^{\mathrm{jet}}_{\mathrm{rad}}(\delta),
      \mathscr N_\xi
    \big\rangle_{\cH_S}
  }
\end{equation}
for every \(\delta\in\Gamma_{\mathrm{rad}}\).  Equivalently,
\begin{equation}\label{eq:dual-news-definition}
  \mathscr N_\xi
  =
  \big((\cN^{\mathrm{jet}}_{\mathrm{rad}})^{-1}\big)^\star
  \mathcal F_\xi^{\mathrm{rad}} .
\end{equation}

The representation has the expected structural properties:
\begin{align}
  \mathscr N_{a\xi+b\eta}
  &=
  a\mathscr N_\xi+b\mathscr N_\eta ,\\
  \cN^{\mathrm{jet}}_{\mathrm{rad}}(\delta)=0
  &\Longrightarrow
  \mathcal F_\xi^{\mathrm{rad}}(\delta)=0\quad\forall\xi,\\
  \big(\forall\delta,\ \mathcal F_\xi^{\mathrm{rad}}(\delta)=0\big)
  &\Longleftrightarrow
  \mathscr N_\xi=0 .
\end{align}
The first line is \(\xi\)-linearity, the second is the unified zero set, and the last follows because \(\cN^{\mathrm{jet}}_{\mathrm{rad}}(\Gamma_{\mathrm{rad}})\) is dense in its closure.  Gauge invariance of \(\mathscr N_\xi\) follows from assumption 3 and invariance of \(\Omega\) and \(X_\xi\).

This is a flux representation, but not a news-square theorem.  The news enters \eqref{eq:dual-news-flux} linearly, while \(\mathscr N_\xi\) carries the remaining symplectic argument and the generator.  A universal square law would require \(\mathscr N_\xi=c_\xi\,\cN^{\mathrm{jet}}_{\mathrm{rad}}\) with \(c_\xi\) independent of phase-space position; the sign-flip obstruction and the disjoint-support charge-blindness example exclude this on the full linearized phase space.  Positivity would require the much stronger restriction that, on a chosen outgoing cone, \(\mathscr N_\xi\) is proportional to the news with a positive coefficient.  No such restriction has been proved.

The representation also clarifies the remaining open step.  The detector zero set and detector magnitude are already encoded in \(\cN^{\mathrm{jet}}_{\mathrm{rad}}\).  What remains is to show that the dual fields \(\mathscr N_\xi\) admit local or conformally covariant representatives on arbitrary cuts, to extend the construction beyond strongly equipped outgoing data, and to prove the controlled flat limit of the pairing.  The finite logical core is formalized by \texttt{SymplecticDualNews.Model}, \texttt{flux\_zero\_of\_news\_zero}, \texttt{flux\_zero\_all\_iff\_dual\_zero}, and \texttt{not\_square\_of\_sign\_flip}.

\subsubsection{Local adjoint representative in the linearized sector}

\statusD The Riesz representative in \eqref{eq:dual-news-definition} can be written explicitly in the linearized three-dimensional boundary theory.  Let \(D_a\) be the covariant derivative of the conformally rescaled boundary metric \(\mathring q_{ab}\) on \(\scri^+\), let \(\beta_\xi\) be the normal component of the spacetime Killing field, and set
\begin{equation}\label{eq:boundary-symmetry-operator}
  A_\xi h_{ab}
  :=
  \mathcal L_\xi h_{ab}+2H\beta_\xi h_{ab}
  =
  \xi^cD_ch_{ab}
  +
  h_{cb}D_a\xi^c
  +
  h_{ac}D_b\xi^c
  +
  2H\beta_\xi h_{ab} .
\end{equation}
Then the Hamiltonian flux is \(\mathcal F_\xi=\varkappa\int E^{ab}A_\xi h_{ab}\,\diff\mu\).  Integrating the first term by parts and discarding the compact-support or Schwartz boundary term gives
\begin{equation}\label{eq:local-dual-news-adjoint}
  \boxed{
    \mathscr N_\xi^{\mathrm{loc}\,ab}
    =
    \varkappa
    \Big[
      -
      D_c\big(\xi^c E^{ab}\big)
      +
      2E^{c(a}D_c\xi^{b)}
      +
      2H\beta_\xi E^{ab}
    \Big]
  }
\end{equation}
and
\begin{equation}\label{eq:local-dual-news-pairing}
  \mathcal F_\xi
  =
  \int_{\scri^+}
  \mathscr N_\xi^{\mathrm{loc}\,ab}h_{ab}\,\diff\mu .
\end{equation}
Indeed, the two terms containing \(D_a\xi^c\) combine to \(2h_{ab}E^{c(a}D_c\xi^{b)}\).  Thus the dual news is not merely an abstract Riesz vector: in the linearized sector it is the local adjoint of the boundary symmetry operator acting on the linearized Weyl electric field.

Let \(\Pi_{\mathrm{rad}}\) be the \(L^2\)-orthogonal projection onto the transverse-traceless radiative tensors.  Since the physical perturbation satisfies \(h_{ab}=\Pi_{\mathrm{rad}}h_{ab}\), the canonical representative in the radiative Hilbert space is
\begin{equation}\label{eq:local-dual-news-radical}
  \mathscr N_{\xi,\mathrm{rad}}^{\mathrm{loc}}
  :=
  \Pi_{\mathrm{rad}}\mathscr N_\xi^{\mathrm{loc}},
  \qquad
  \mathcal F_\xi
  =
  \left\langle
    \mathscr N_{\xi,\mathrm{rad}}^{\mathrm{loc}},h
  \right\rangle_{L^2} .
\end{equation}
Uniqueness of the orthogonal projection gives
\(\mathscr N_{\xi,\mathrm{rad}}^{\mathrm{loc}}=\mathscr N_\xi\) in the notation of \eqref{eq:dual-news-flux}, whenever the hypotheses of that representation hold.  The formula is manifestly \(\xi\)-linear because \(A_{a\xi+b\eta}=aA_\xi+bA_\eta\), and it is gauge invariant after projection because \(E_{ab}\) is gauge invariant and \(\Pi_{\mathrm{rad}}\) annihilates pure gauge directions.  For the six boundary-tangent Euclidean generators \(\beta_\xi=0\); for the de Sitter time translation \(\beta_\xi=1\), so the extra conformal-mass term is exactly the term responsible for the nonzero energy flux.  In the controlled no-incoming sector \(E_{ab}\) is fixed, up to the Hodge rotation, by the tensor detector \(\hat C_{AB}\); hence \eqref{eq:local-dual-news-adjoint} is a local representative built directly from the detector and the generator.

This is still not a news-square theorem: \eqref{eq:local-dual-news-pairing} pairs \(E\)-data with \(h\)-data.  A universal pointwise square law would require a dynamical relation identifying \(E_{ab}\) with a positive multiple of \(A_\xi h_{ab}\) on the full radiative phase space, i.e. precisely the additional incoming/outgoing constraint that fails there.  The finite linearity algebra for the local adjoint is formalized by \texttt{LocalDualNews.localDual}, \texttt{localDual\_add}, \texttt{localDual\_smul}, \texttt{PairingAlgebra}, and \texttt{pairing}.

\paragraph{Retarded adjoint factorization criterion.}
\statusD The exact condition closing the remaining linear gap can be stated without choosing a source model.  Let \(P_{\mathrm{rad}}\) denote radiative TT projection and let \(\langle\cdot,\cdot\rangle_\kappa\) denote the inner product with weight \(\varkappa\) that appears in \(\mathcal F_\xi=\varkappa\int E^{ab}A_\xi h_{ab}\,\diff\mu\).  For a fixed generator \(\xi\), suppose that the retarded sector admits a news seed \(\nu_\xi\) and linear operators \(B_\xi,C_\xi\) on radiative fields such that
\begin{equation}\label{eq:retarded-factorization}
  P_{\mathrm{rad}}A_\xi h
  =
  B_\xi\nu_\xi ,
  \qquad
  P_{\mathrm{rad}}E
  =
  C_\xi\nu_\xi ,
  \qquad
  \varkappa\big\langle C_\xi\nu,B_\xi\nu\big\rangle
  =
  c_\xi\|\nu\|^2
\end{equation}
for some phase-space-independent constant \(c_\xi\ge0\).  Then
\begin{equation}\label{eq:retarded-square-law}
  \mathcal F_\xi
  =
  c_\xi\,
  \|\nu_\xi\|^2 ,
  \qquad
  n_\xi:=\sqrt{c_\xi}\,\nu_\xi ,
\end{equation}
so \(\mathcal F_\xi=\|n_\xi\|^2\) and \(\mathcal F_\xi=0\iff n_\xi=0\) when \(c_\xi>0\).  Equivalently, \(\varkappa C_\xi^\dagger B_\xi=c_\xi I\); an integration by parts is precisely what supplies this adjoint identity.  The stronger pointwise requirement \(P_{\mathrm{rad}}E=c_\xi P_{\mathrm{rad}}A_\xi h\) is a special sufficient case, obtained when \(C_\xi\) is a multiple of \(B_\xi\), but it is not necessary.  The finite algebra is formalized by \texttt{RetardedSquareCriterion.Model}, \texttt{square}, \texttt{nonnegative}, and \texttt{zero\_iff}.

\statusL The adjoint factorization is realized in the retarded first-post-de Sitter quadrupole sector of Ashtekar--Bonga--Kesavan \cite{ashtekar2015quadrupole}.  There one has
\(\cE_{ab}=(2H)^{-1}(1/\eta)(\partial_\eta^2-\eta^{-1}\partial_\eta)\chi^{\mathrm{TT}}_{ab}\) and
\(A_T\chi^{\mathrm{TT}}_{ab}\hat{=}T^m\mathring\nabla_m\chi^{\mathrm{TT}}_{ab}\).  With the retarded solution, substitution in the energy flux gives two radial-source factors which are not pointwise proportional; after integration by parts their pairing is the square of the source radiation field
\(\cR^{\mathrm{src}}_{ab}:=\partial_r\!\int\!\diff^3x'\,T^{\mathrm{TT}\,'}_{ab}(\eta_r,\vec x')\).  Thus
\begin{equation}\label{eq:quadrupole-source-square}
  E_T
  =
  \frac{G}{2\pi}
  \int_{\scri^+}
  \diff T\,\diff^2S\,
  \cR^{\mathrm{src}}_{ab}\cR_{\mathrm{src}}^{ab}
  \ge0 .
\end{equation}
This is the first controlled linear sector in which the local dual-news pairing becomes an actual positive square; the news seed is the source radiation field itself rather than a pointwise multiple of \(A_T\chi^{\mathrm{TT}}\).  It does not contradict the sign-flip obstruction: the factorization \eqref{eq:retarded-factorization} is a retarded-source constraint and fails on the full homogeneous de Sitter phase space.  It also does not extend automatically to the nonlinear AldS charge flux or to the FAS first-jet detector square.

\subsection{Conformal phase-space flux}

\statusL Hoque--Krtou{\v{s}}--Pe{\'o}n-Nieto show that the flat-space Geroch--Xanthopoulos regularization fails for \(\Lambda>0\): flat perturbations can be rescaled to \(O(\Omega)\), while \(\Lambda>0\) consistency forces \(O(1)\) falloff \cite{hoque2025conformal}.  Nevertheless the presymplectic flux is finite and nonzero:
\begin{equation}
  \bm F_\xi
  =
  \frac{\ell}{16\pi G}
  \int_{\scri^+}
  |\tilde\gamma|^{1/2}
  \delta\tilde e^{ab}
  \cL_{\tilde\xi}\delta\tilde\gamma_{ab}.
\end{equation}
This is a phase-space flux at spacelike \(\scri^+\); it does not identify a local news tensor or a mass-aspect positivity theorem.

\section{Exact-solution benchmark}

The detector has been tested on the available exact families.  The following table summarizes the outcomes; it validates the detector, not a news-square theorem.

\begin{center}
\begin{tabular}{p{0.27\linewidth}p{0.64\linewidth}}
\hline
Solution & Detector / flux status\\
\hline
de Sitter & \(C_{ab}=\cF_{ab}=0\), \(p_a=0\); exact nonradiating reference.\\[1mm]
Kottler / Schwarzschild--de Sitter & Spherical symmetry gives commuting Cotton and stress; \(p_a=0\).  AldS charges are nonzero mass charges.\\[1mm]
Kerr--de Sitter & Cotton and stress are proportional through the induced conformal Killing field; \(p_a=0\).\\[1mm]
Robinson--Trautman-dS & Generic boundary Cotton and stress fail to commute, so \(p_a\neq0\).  Restricted monotonic charges exist; positivity is representative-dependent.\\[1mm]
Photon rockets & In an oriented frame,
\(\bar{\cP}^{\hat a}=-2(p^2+q^2)\delta^{\hat a}{}_{\hat1}+4q(m-p)\delta^{\hat a}{}_{\hat2}\);
vanishing iff \(p=q=0\), the generalized Kinnersley cases \cite{fabSen2026rockets}.\\[1mm]
C-metric & \(\bar{\cP}^a\propto\alpha m^2\delta^a_p\), nonzero for acceleration \(\alpha\neq0\) \cite{fabSen2020radiation}.\\[1mm]
\hline
\end{tabular}
\end{center}

For Robinson--Trautman-dS, \cite{poole2025charges} also finds monotonic charges for suitable representatives and Calabi functionals.  With the Bondi-time translation and a different representative, the flux contains a \(\Lambda\)-dependent term
\begin{equation}
  \frac{\diff E}{\diff u_B}
  =
  -\frac{\epsilon^2}{16\pi G}
  \oint \diff^2 S
  \left[
    (\dot f''-\cot\theta\,\dot f')^2
    +\frac{3m}{\ell^2}\partial_{u_B}(f'^2)
  \right],
\end{equation}
so positivity is not representative-independent even in this restricted class.

\section{The \(\Lambda\to0\) limit}

\statusL In a controlled retarded one-parameter family with source scale \(D_o\ll\ell\), \(H\Delta t_{\rm source}\ll1\), and finite source duration, Ashtekar--Bonga--Kesavan show that the de Sitter quadrupole flux approaches the Einstein quadrupole formula in the physical \((t,\vec x)\) chart \cite{ashtekar2015quadrupole}.  The tail correction is \(O(H)\) and vanishes as \(H\to0\).  The flux has the schematic form
\begin{equation}
  E_T
  =
  \frac{G}{8\pi}
  \int_{\scri^+}\diff T\,\diff^2S\,
  \cR_{ab}\cR^{TT}_{cd}\mathring q^{ac}\mathring q^{bd},
\end{equation}
with \(\cR_{ab}\) reducing, in the nonrelativistic short-duration regime, to the third time derivative of the mass quadrupole plus \(O(H)\) corrections.

\subsubsection{Flat flux bridge from the retarded source news seed}
\statusD The preceding flux law can be sharpened to an explicit flat-limit news normalization.  Let \(u=\eta_r=\eta-r\) be the retarded source time and let
\(\nu^{\mathrm{src}}_{ab}:=\partial_r\int \diff^3x'\,T^{\mathrm{TT}\, '}_{ab}(u,\vec x')\).  Conservation of the leading Newtonian source gives \(\int T_{ij}\diff^3x=\tfrac12\ddot Q^{(\rho)}_{ij}\), so at \(H=0\)
\begin{equation}\label{eq:source-seed-quadrupole}
  \nu^{\mathrm{src}}_{ab}
  =
  -\tfrac12
  P^{\mathrm{TT}}\dddot Q^{(\rho)}_{ab}+O(H) .
\end{equation}
On the other hand, the standard quadrupole strain in the far zone is
\(h^{\mathrm{TT}}_{AB}=(2G/r)\ddot Q^{\mathrm{TT}}_{AB}(u)\), hence the Bondi news is
\begin{equation}\label{eq:flat-bondi-news-from-seed}
  N^{\mathrm{Bondi}}_{AB}
  =
  \partial_u\!\left(rh^{\mathrm{TT}}_{AB}\right)
  =
  \frac{2G}{r}P^{\mathrm{tt}}\dddot Q^{(\rho)}_{AB}
  =
  -\frac{4G}{r}\,
  P^{\mathrm{tt}}\nu^{\mathrm{src}}_{AB}
  +O(H) .
\end{equation}
Consequently the flat Bondi mass-loss density and the ABK source-square density have the same normalization after the observer's radial \(\mathrm{tt}\)-projection:
\begin{equation}\label{eq:flat-flux-bridge}
  \frac{r^2}{32\pi G}
  N^{\mathrm{Bondi}}_{AB}N_{\mathrm{Bondi}}^{AB}
  =
  \frac{G}{2\pi}
  \big(P^{\mathrm{tt}}\nu^{\mathrm{src}}\big)_{AB}
  \big(P^{\mathrm{tt}}\nu^{\mathrm{src}}\big)^{AB}.
\end{equation}
Integrating over \(u\) and \(S^2\) gives
\[
  \lim_{H\to0}E_T^{\mathrm{ABK}}
  =
  \frac{G}{8\pi}
  \int \diff u\,\diff^2S\,
  \dddot Q^{\mathrm{TT}}_{ab}\dddot Q_{\mathrm{TT}}^{ab}
  =
  \frac{1}{32\pi G}
  \int \diff u\,\diff^2S\,
  N^{\mathrm{Bondi}}_{AB}N_{\mathrm{Bondi}}^{AB}.
\]
The first equality is the ABK flat limit, independently confirmed to leading order by the comparison of the three quadrupole derivations \cite{ashtekar2015quadrupole,dobkowski2025comparison}; the Bondi flux normalization is standard \cite{wald1984general}.  Thus the retarded source news seed has a controlled flat-limit representative \(P^{\mathrm{tt}}\nu^{\mathrm{src}}\), and \(-(4G/r)^{-1}N^{\mathrm{Bondi}}\) is its tensorial Bondi image.  This does not identify the full three-dimensional source field with ordinary Bondi news: only the radial \(\mathrm{tt}\)-projected part enters the Bondi flux, and no analogous statement has been proved for generic \(\Lambda>0\) data.  The finite normalization algebra is formalized by \texttt{FlatBondiBridge.Model}, \texttt{flux\_equivalence}, and \texttt{scaled\_seed\_square}.

\statusL Fernández--Álvarez and Senovilla give a separate flat-limit constraint at the level of the super-Poynting detector \cite{fabSen2022asymptotic}.  Assume that \(\lim_{\Lambda\to0}g_{\alpha\beta}\) exists and defines a good Lorentzian metric.  Their canonical asymptotic supermomentum obeys
\begin{equation}
  \lim_{\Lambda\to0}p^\alpha
  =
  Q^\alpha_{\mathrm{flat}}
  =
  -\big(N^\mu N^\nu N^\rho \mathcal D_{\alpha\mu\nu\rho}\big)\big|_{\Lambda=0}.
\end{equation}
Hence criterion 1 in the \(\Lambda>0\) theory, namely \(p^\alpha=0\), forces \(Q^\alpha_{\mathrm{flat}}=0\) whenever the regular limit exists; in the flat theory this is equivalent to vanishing of the Bondi news.  For the radiant supermomenta \(q^\pm\) obtained from the rescaled null fields with \(\|M^\alpha\|^2\propto\Lambda^p\), \(p>0\), their Lemma IV.6 gives
\begin{equation}
  \lim_{\Lambda\to0}q^\pm{}^\alpha
  =
  \frac{(1\pm B)^3}{2\sqrt2}\,Q^\alpha_{\mathrm{flat}},
  \qquad
  \lim_{\Lambda\to0}q^+\!=0=\lim_{\Lambda\to0}q^-
  \iff
  Q^\alpha_{\mathrm{flat}}=0 .
\end{equation}
This is a regular-limit implication from the de Sitter detector to the flat detector.  It does not show that \(V_{AB}\) or \(n^+_{AB}=2V_{AB}\) has a flat limit, nor that the converse implication holds without the regularity hypothesis.

\statusD The \(q^\pm\) relation gives a sharper statement for the controlled outgoing sector.  There the no-incoming condition is \(Q_-^\mu=0\) for every \(\Lambda>0\), hence the rescaled incoming radiant momentum \(q^-_\alpha\) vanishes identically.  Taking the limit in the FAS formula gives
\begin{equation}\label{eq:flat-no-incoming-alternative}
  0
  =
  \lim_{\Lambda\to0}q^-{}^\alpha
  =
  \frac{(1-B)^3}{2\sqrt2}\,Q^\alpha_{\mathrm{flat}} .
\end{equation}
Therefore
\begin{equation}\label{eq:flat-degeneracy-dichotomy}
  Q^\alpha_{\mathrm{flat}}\neq0
  \quad\Longrightarrow\quad
  B=1 .
\end{equation}
Equivalently, a regular flat limit with nonzero flat Bondi news from a no-incoming de Sitter family must be directionally degenerate: \(M^\alpha=N m^\alpha\to N^\alpha|_{\Lambda=0}\), so the incoming cut-null direction \(K^-_\alpha=(N_\alpha-M_\alpha)/\sqrt2\) collapses, \(\lim K^-=0\) by FAS eq. (IV.93).  If instead \(B\neq1\), no incoming radiation at \(\scri^+\) forces \(Q^\alpha_{\mathrm{flat}}=0\), hence vanishing flat Bondi news.  The same argument with \(q^+\) shows that no outgoing radiation at \(\scri^+\) forces either \(B=-1\) or \(Q^\alpha_{\mathrm{flat}}=0\).

This dichotomy explains why a flat limit of the first-jet norm is not automatic.  The identity \(\|\cN^{\mathrm{jet}}_+\|^2=-\tfrac12\int p_m\,dS\) uses the normalized cut normal \(m^a=N^{-1}M^a\), which becomes singular as \(N=\sqrt{\Lambda/3}\to0\).  In the nondegenerate case \(B\neq\pm1\) the flat detector already vanishes.  In the radiating outgoing case one needs the degenerate regime \(B=1\), and a finite limit of \(p_m\) requires additional control of the singular cut direction.  The scaling of the controlled first-jet object can now be made explicit.

\subsubsection{Exact first-jet scaling in the flat limit}

\statusD Decompose the rescaled outgoing supermomentum in the finite \(K^\pm=Nk^\pm\) basis,
\begin{equation}\label{eq:qK-basis}
  q^+_\alpha
  =
  \alpha_+ K^+_\alpha
  +
  \beta_+ K^-_\alpha
  +
  \gamma_+{}^A E^+_{A\alpha} .
\end{equation}
Comparing this with the FAS decomposition
\(Q^+=W^+k^++Z^+k^-+\cdots\) and using \(q^+=N^3Q^+\), \(K^\pm=Nk^\pm\), gives
\begin{equation}\label{eq:alpha-beta-W-Z}
  \alpha_+=N^2W^+,
  \qquad
  \beta_+=N^2Z^+ .
\end{equation}
In the controlled outgoing sector the no-incoming identifications give
\(\hat C^{(+)}_{AB}=2\hat C_{AB}\) and \(C^{(+)}_A=C_A\), hence
\begin{equation}
  W^+=4\hat C_{AB}\hat C^{AB},
  \qquad
  Z^+=4C_AC^A .
\end{equation}
Using \(\|\Phi^{\mathrm{norm}}\|^2=\hat C^2=W^+/4\) and
\(\|C_A[n^+]\|^2=Z^+/4\), equations \eqref{eq:alpha-beta-W-Z} yield the exact pointwise scaling laws
\begin{equation}\label{eq:first-jet-flat-scaling}
  N^2\|\Phi^{\mathrm{norm}}\|^2
  =
  \frac{\alpha_+}{4},
  \qquad
  N^2\|C_A[n^+]\|^2
  =
  \frac{\beta_+}{4},
  \qquad
  N^2\|\cN^{\mathrm{jet}}_{\delta,\gamma}\|^2
  =
  \frac{\alpha_+}{4}+\frac{\beta_+}{4}.
\end{equation}
The integrated form follows under the usual dominated-convergence hypothesis.

This separates two regimes.  First, suppose \(B=1\), the only nondegenerate source of flat radiation in the no-incoming sector.  Then \(K^+\to\sqrt2\,N^\alpha|_{\Lambda=0}\) is finite and nonzero, while \(\lim q^+=2\sqrt2\,Q_{\mathrm{flat}}\).  Hence \(\alpha_+\) has a finite limit \(\alpha_\infty\), nonzero iff \(Q_{\mathrm{flat}}\neq0\), and
\begin{equation}\label{eq:tensor-jet-flat-limit}
  \lim_{\Lambda\to0}
  N^2\|\Phi^{\mathrm{norm}}\|^2
  =
  \frac{\alpha_\infty}{4}.
\end{equation}
Thus the tensor component of the first-jet news renormalizes universally by \(N\): \(N\Phi^{\mathrm{norm}}\) has a finite nonzero norm exactly when the flat Bondi news is nonzero.  This is a genuine flat-limit representative of the tensor detector sector.

Second, the vector component is controlled by \(\beta_+\).  If \(B\neq\pm1\), both \(K^+\) and \(K^-\) have finite nonzero limits, so \(\beta_+\) is finite; the no-incoming condition then forces \(Q_{\mathrm{flat}}=0\), and \eqref{eq:first-jet-flat-scaling} gives zero flat limit for the whole first jet.  If \(B=1\), however, \(K^-\to0\); finiteness of \(q^+\) alone permits \(\beta_+\) to diverge at a rate determined by the collapse of \(K^-\).  Therefore the flat limit of \(C_A[n^+]\), and hence of the full first-jet news and of \(V_{AB}\) and \(n^+_{AB}\), is not determined by the FAS rescaled momentum limit.  The missing datum is precisely the degeneracy rate of \(K^-\), equivalently the limiting coefficient \(\beta_+\).

\paragraph{Bounded-detector resolution of \(\beta_+\).}
\statusD The exact identity \(\beta_+=4N^2\|C_A[n^+]\|^2\) makes the nature of this datum concrete.  If the vector detector stays bounded in a fixed conformal frame as \(\Lambda\to0\), then automatically \(\beta_+\to0\), and the whole \(N\)-renormalized first jet has the same flat norm as its tensor component:
\begin{equation}\label{eq:bounded-detector-beta}
  C_A[n^+]=O(1)
  \quad\Longrightarrow\quad
  \beta_+=O(N^2)\to0,
  \qquad
  \lim_{\Lambda\to0}N^2\|\cN^{\mathrm{jet}}_+\|^2
  =
  \frac{\alpha_\infty}{4}.
\end{equation}
Conversely, a finite nonzero \(\beta_\infty\) requires the magnetic detector component to diverge exactly as \(N^{-1}\) in the collapsing cut-normal frame, while a divergent \(\beta_+\) requires faster growth.  Thus \(\beta_+\) is not an additional independent radiation coefficient: it is precisely the growth rate of the vector detector along the incoming null direction that collapses when \(B=1\).  In particular, no extra datum is needed for any controlled retarded family in which \(C_A[n^+]\) is uniformly bounded.  Whether that boundedness follows from the ABK retarded quadrupole fields remains open because the normalized cut normal itself becomes singular.  The finite algebra is formalized by \texttt{BetaPlusBound.Model}, \texttt{beta\_eq}, \texttt{beta\_nonneg}, \texttt{beta\_le\_of\_bound}, and \texttt{beta\_zero\_of\_N\_zero}.  The original scaling algebra is formalized by \texttt{FirstJetFlatScaling.Algebra}, \texttt{scaled\_phi\_eq}, \texttt{scaled\_vector\_eq}, \texttt{scaled\_jet\_eq}, and the corresponding \texttt{tendsto} lemmas.

\statusD The sign-flip operation in the sign-flip obstruction can be made explicit in the TT family of \cite{ashtekar2015linear}.  Writing
\begin{equation}
  E_{\vec k}^{(s)}=c_{\vec k}^{(s)}\sin(k/H)+d_{\vec k}^{(s)}\cos(k/H),
  \qquad
  B_{\vec k}^{(s)}=-c_{\vec k}^{(s)}\cos(k/H)+d_{\vec k}^{(s)}\sin(k/H),
\end{equation}
the transformation \(E\mapsto-E\), \(B\mapsto B\) acts on \((c,d)\) by
\begin{equation}\label{eq:sign-flip-flat}
  \begin{pmatrix}c'\\ d'\end{pmatrix}
  =
  \begin{pmatrix}
    \cos(2k/H) & -\sin(2k/H)\\
    -\sin(2k/H) & -\cos(2k/H)
  \end{pmatrix}
  \begin{pmatrix}c\\ d\end{pmatrix}.
\end{equation}
For coefficients \(c,d\) with a smooth \(H\to0\) limit, the transformed coefficients generically oscillate infinitely fast.  Thus the sign reversal that forbids a sign-definite square on the full de Sitter phase space typically leaves the controlled flat-limit family.  This reconciles the sign-flip obstruction with the positive-definite flat energy without asserting a universal \(\Lambda\to0\) theorem.

This is a controlled limit of special families, not a universal boundary limit.  Reasons include:
\begin{itemize}
  \item linearized dS energy can be arbitrarily negative for every \(\Lambda>0\), whereas flat energy is positive definite \cite{ashtekar2015linear};
  \item all dS Killing fields become spacelike near \(\scri^+\), so the notion of time translation degenerates;
  \item flux formulas can pass continuously to Bondi form while the asymptotic symmetry algebra jumps from the finite \(\Lambda>0\) algebra to BMS \cite{bonga2023radiation};
  \item recovery of no-incoming Bondi data has subtle differences \cite{ashtekar2019noincoming};
  \item the Bondi-gauge relation \eqref{eq:gamma1-obstruction} requires the \(O(\Lambda)\) scaling condition for regularity.
\end{itemize}

A candidate AldS news must therefore come with both an outgoing/no-incoming condition and a scaling prescription.  The detector alone does not provide either.

\section{Lean formalization of the conditional core}

Two Lean files were produced and verified.

\subsection{Core scaffold}

\texttt{LambdaPositiveBondi\_Plan.lean} uses integers and no Mathlib import.  It proves the abstract flux lemmas
\texttt{conserved\_of\_flux\_free}, \texttt{charge\_nonincreasing}, \texttt{equal\_charge\_iff\_zero\_flux}, and a toy model \texttt{toy\_charge\_nonincreasing}.  The open physics is represented by structures \texttt{RigidityTarget} and \texttt{DesiredNewsConstruction}, not by axioms.

\subsection{Mathlib scaffold}

\texttt{pnp-mathlib/PNP/LambdaPositiveBondi.lean} uses real-valued Bochner integrals and contains:
\begin{itemize}
  \item \texttt{AldSFluxModel}, \texttt{FluxFreeVector}, and the associated flux/logic lemmas;
  \item finite-dimensional boundary algebra: \texttt{SymTensor3}, \texttt{Cotton}, \texttt{Stress}, \texttt{commutator}, \texttt{superPoynting};
  \item the \(\Lambda>0\) equivalence \texttt{canonicalSuperPoyntingZero\_iff\_commutatorZero} and \texttt{noRadiation\_iff\_commuting};
  \item the \(\Lambda<0\) proportionality target and the explicit commuting-but-not-proportional example;
  \item exact-solution algebra \texttt{rocket\_noRadiation\_iff} and \texttt{outgoing\_noRadiation\_iff};
  \item outgoing flux positivity \texttt{OutgoingCutData.normalFlux\_nonpos};
  \item the conditional cut-news skeleton \texttt{CutNews.CutNewsCandidate}, with \texttt{source\_zero\_iff\_news\_zero}, \texttt{news\_zero\_iff\_detector\_zero}, and the contrapositive \texttt{news\_nonzero\_iff\_radiation};
  \item the linearized algebra \texttt{scaleSym}, \texttt{commutator\_scaleSym\_scaleSym}, \texttt{superPoynting\_scaleSym\_scaleSym}, and \texttt{linearized\_superPoynting\_is\_quadratic};
  \item the flux-linearity obstruction \texttt{LinearFluxObstruction.no\_nonzero\_constant\_square};
  \item the linearized Hamiltonian news-flux core \texttt{HamiltonianNewsFlux.Model}, \texttt{flux\_smul}, \texttt{flux\_add}, and \texttt{flux\_zero};
  \item the conformal-Killing degeneracy core \texttt{CKVFluxDegeneracy.trace}, \texttt{ckvDerivative}, and \texttt{contract\_zero};
  \item the pointwise news-square obstruction \texttt{NewsSquareObstruction.no\_linear\_pointwise\_square} and the harmonic-mode obstruction \texttt{HarmonicDimensionObstruction.no\_injective\_pair\_to\_tensor};
  \item the linearized sign-flip obstruction \texttt{LinearizedFluxObstruction.SignFlipNewsModel}, \texttt{energy\_zero\_of\_sign\_flip\_square}, and \texttt{no\_sign\_flip\_square};
  \item the conditional strongly equipped radiant-news core \texttt{RadiantCutNews.StrongModel}, \texttt{nPlus\_zero\_iff\_V\_zero}, \texttt{Zplus\_iff\_V\_zero}, and \texttt{V\_zero\_iff\_Pzero};
  \item the inverse-divergence core \texttt{RadiantCutNews.StrongOutgoingInversion.Model}, \texttt{div\_V\_eq\_half\_source}, and \texttt{nPlus\_eq\_twoV\_of\_div};
  \item the detector-square representative \texttt{DetectorSquareNews.Pair}, \texttt{DetectorSquareNews.Model}, and \texttt{DetectorSquareNews.normalFlux\_zero\_iff\_news\_zero};
  \item the minimal rank-two tensor-pair core \texttt{RankTwoTensorPairNews.VectorTensorIsometry} and \texttt{normalFlux\_zero\_iff\_news\_zero};
  \item the radiant transport-square core \texttt{RadiantNewsTransport.Model}, \texttt{W\_nonneg}, and \texttt{W\_zero\_iff\_nDot\_zero};
  \item the outgoing zero-set unification core \texttt{NewsZeroSetUnification.Model}, \texttt{detector\_zero\_iff\_radiant\_zero}, and \texttt{charge\_iff\_noRadiation};
  \item the spectral magnitude obstruction \texttt{SpectralNormMismatch.Model} and \texttt{no\_uniform\_lower\_bound};
  \item the first-jet radiant-news core \texttt{RadiantJetNews.Model}, \texttt{jet\_sq\_eq\_neg\_half\_normalFlux}, and \texttt{jet\_zero\_iff\_detector\_zero}, together with the news-only square algebra \texttt{NewsOnlyJetSquare.Algebra} and \texttt{scaled\_jet\_sq};
  \item the generalized \((\delta,\gamma)\) inversion core \texttt{GeneralizedJetRotation.TensorRotation}, \texttt{sq\_inverse}, and \texttt{inverse\_zero\_iff};
  \item the charge-blindness core \texttt{ChargeFluxBlindness.Model}, \texttt{news\_zero\_flux\_zero}, and \texttt{not\_flux\_zero\_imp\_news\_zero};
  \item the conformal-weight core \texttt{ConformalWeightObstruction.LocalTensorOperator}, \texttt{tensor\_square\_weight}, \texttt{vector\_square\_weight}, and \texttt{no\_local\_rank\_two\_map};
  \item the arbitrary-cut polar-decomposition core \texttt{ArbitraryCutJetNews.PolarData}, \texttt{toIsometry}, and \texttt{J\_zero\_iff};
  \item the conformal-intertwiner weight algebra \texttt{ConformalJetIntertwiner.Algebra}, \texttt{input\_weight}, \texttt{output\_weight}, and \texttt{jet\_square\_scaling}, together with \texttt{PolarAlgebra} and \texttt{polar\_unique};
  \item the symplectic dual-news core \texttt{SymplecticDualNews.Model}, \texttt{flux\_zero\_of\_news\_zero}, and \texttt{flux\_zero\_all\_iff\_dual\_zero}, together with the local-adjoint algebra \texttt{LocalDualNews.localDual}, \texttt{localDual\_add}, \texttt{localDual\_smul}, \texttt{PairingAlgebra}, \texttt{pairing}, the retarded factorization core \texttt{RetardedSquareCriterion.Model}, \texttt{square}, \texttt{nonnegative}, \texttt{zero\_iff}, and the flat flux normalization \texttt{FlatBondiBridge.Model}, \texttt{scaled\_seed\_square}, \texttt{flux\_equivalence}, together with the bounded-detector estimate \texttt{BetaPlusBound.Model}, \texttt{beta\_eq}, \texttt{beta\_nonneg}, \texttt{beta\_le\_of\_bound};
  \item the flat-limit directional degeneracy core \texttt{FlatLimitDirection.Model} and \texttt{qMinus\_zero\_Qflat\_ne\_zero\_imp\_B\_one};
  \item the first-jet flat-scaling core \texttt{FirstJetFlatScaling.Algebra}, \texttt{scaled\_phi\_eq}, \texttt{scaled\_vector\_eq}, and \texttt{scaled\_jet\_eq};
  \item the regular flat-limit implication \texttt{FlatLimitConstraint.Model} and \texttt{flat\_news\_zero}.
\end{itemize}

Both files pass:
\begin{verbatim}
lean LambdaPositiveBondi_Plan.lean
cd pnp-mathlib
lake env lean PNP/LambdaPositiveBondi.lean
lake build PNP.LambdaPositiveBondi
\end{verbatim}

The formalization deliberately stops short of asserting the unresolved geometric theorem.

\section{Conclusion and open problems}

The main open problem can now be stated more precisely than in the usual informal form.

\begin{problem}[\statusO: refined news problem]
For admissible AldS boundary data \((h_{ij},\cT_{ij})\), construct a gauge-covariant cut tensor or bilinear news object \(\mathcal N\) such that:
\begin{enumerate}
  \item its vanishing is equivalent to the radiation detector \([C,\cF]=0\), at least in the outgoing sector;
  \item its pairings with the asymptotic translation charges reproduce the AldS flux law \eqref{eq:aldS-flux}, respecting the linearity in \(\xi\);
  \item the flux is nonpositive and equality is characterized by no radiation;
  \item it reduces, under a controlled scaling regime, to the Bondi news \(N^{\mathrm{Bondi}}_{AB}\) as \(\Lambda\to0\);
  \item it admits a scattering interpretation with an incoming/outgoing split.
\end{enumerate}
\end{problem}

Specific remaining tasks are:
\begin{enumerate}
  \item Fix precise function spaces and conformal rescaling conventions for the AldS state space.
  \item Extend the strongly equipped identity \(n^+_{AB}=2V_{AB}\) to a general \(X_A\).  Equivalently, construct a weight-\(\omega^{-2}\), CKV-orthogonal one-form \(X_A\), or a gauge-invariant function \(f\), whose constancy matches no radiation.
  \item Extend \(\cN^{\mathrm{jet}}_+\) beyond strongly equipped outgoing data and construct the conformally covariant analogue of \(\mathcal J\) on arbitrary cuts.  By the weight-parity obstruction, this must be a nonlocal operator of weight \(+1\), not a local curvature correction to \(\mathcal K_q\).
  \item Prove a flux representation for \(\cN^{\mathrm{jet}}_+\), first for linearized de Sitter waves and then generally, using fixed-\(\xi\), \(\xi\)-weighted, or bilinear AldS/symplectic pairings that are not exhausted by the seven finite charge evaluations.
  \item Prove positivity and rigidity, and determine whether a nonlocal or \(\xi\)-dependent news object is necessary.
  \item Establish a universal no-incoming condition or prove a precise substitute.
  \item Formulate and prove an analogue of peeling, or show that no such analogue exists.
  \item Prove a controlled \(\Lambda\to0\) limit with explicit scaling assumptions.
\end{enumerate}

The positive result in this manuscript is conditional but concrete: under strong equipment, \(S^2\) leaves, and no incoming radiation, the rank-two radiant news is \(n^+_{AB}=2V_{AB}\).  The transport law \eqref{eq:radiant-flux-square} represents the radiant superenergy as the square of the equipment derivative of \(n^\pm\), while the detector-square pair \eqref{eq:detector-news-pair} represents the normal outgoing flux exactly in the outgoing sector, and its round-cut rank-two tensor-pair realization \eqref{eq:rank2-pair-news} saturates the harmonic multiplicity obstruction.  Passing to the first jet in \eqref{eq:first-jet-news} combines these structures: \(\cN^{\mathrm{jet}}_+\) has the unified zero set \eqref{eq:first-jet-zero}, its first component is the normalized FAS flux density, and its total norm is exactly \(-\tfrac12p_m\), which FAS VI.123 rewrites as a local quadratic form in the equipment derivative and divergence of \(n^+\).  The negative results narrow what remains: a single \(\xi\)-independent positive scalar cannot represent all charge fluxes; pointwise and representation-theoretic linear squares of the full outgoing detector pair into one STF tensor cannot reproduce the known positive density; \(n^+=2V\) has a spectral norm mismatch with the detector; the full linearized de Sitter energy flux cannot be a sign-definite news square; all seven de Sitter charge fluxes can vanish on nonzero news, so finite charge evaluations are not a news detector; the conformal weights of the first-jet components force any rank-two realization to use a nonlocal operator of weight \(+1\), excluding local differential constructions; and the naive Bondi shear coefficient is not free when \(\Lambda\neq0\).  The round-cut inversions \eqref{eq:codazzi-inversion} and \eqref{eq:V-round-inversion} show that both the unrestricted cut-Codazzi candidate and the strongly equipped radiant news are nonlocal on each cut.  The arbitrary-cut polar decomposition \eqref{eq:arbitrary-cut-J} extends the vector-to-tensor isometry beyond round cuts, with conformal-Killing modes as the intrinsic replacement for the discarded \(l=0,1\) harmonics; its unweighted version is metric dependent, but the weighted construction \eqref{eq:covariant-J} supplies a conformally covariant weight-\(+1\) intertwiner on arbitrary \(S^2\) cuts, after exclusion of conformal-Killing components, and \eqref{eq:polar-factor-J} identifies it as the unique isometric polar factor of the local operator \eqref{eq:weighted-local-K}.  The flat-limit analysis adds one positive result and one sharp obstruction: the tensor first-jet component renormalizes by \(N=\sqrt{\Lambda/3}\) with limiting square \(\alpha_\infty/4\), where \(\alpha_\infty\) is the outgoing rescaled-supermomentum coefficient along the limiting generator; the vector component and the full first jet require the limiting coefficient \(\beta_+\), which is not fixed when the incoming direction collapses.  The local dual news \eqref{eq:local-dual-news-adjoint} makes the linear flux pairing explicit, and the retarded adjoint factorization \eqref{eq:retarded-factorization} isolates exactly what is needed for a positive square.  That criterion is realized in the retarded quadrupole sector by the positive source-radiation square \eqref{eq:quadrupole-source-square}, so the linear retarded obstruction is not merely heuristic.  Moreover \eqref{eq:flat-flux-bridge} normalizes the radial \(\mathrm{tt}\)-projected source seed so that the ABK flux and the ordinary Bondi mass-loss flux agree exactly in the controlled flat limit.  The decisive open step is no longer a zero-set, detector-magnitude, cut-covariance, linear retarded-square, controlled flat-flux-normalization, or formal \(\beta_+\)-identification compatibility: it is to extend the first-jet construction beyond strongly equipped outgoing data, prove or disprove boundedness of the vector detector in the collapsing retarded frame, pass from the explicit local linearized dual news to a nonlinear and cut-covariant representative, and couple \(\cN^{\mathrm{jet}}_+\) to the nonlinear AldS/symplectic charge-flux law while preserving the retarded factorization.

\bibliographystyle{plainnat}
\bibliography{lambda_positive_bondi}

\end{document}
